% Options for packages loaded elsewhere
\PassOptionsToPackage{unicode}{hyperref}
\PassOptionsToPackage{hyphens}{url}
%
\documentclass[
]{article}
\usepackage{amsmath,amssymb}
\usepackage{iftex}
\ifPDFTeX
  \usepackage[T1]{fontenc}
  \usepackage[utf8]{inputenc}
  \usepackage{textcomp} % provide euro and other symbols
\else % if luatex or xetex
  \usepackage{unicode-math} % this also loads fontspec
  \defaultfontfeatures{Scale=MatchLowercase}
  \defaultfontfeatures[\rmfamily]{Ligatures=TeX,Scale=1}
\fi
\usepackage{lmodern}
\ifPDFTeX\else
  % xetex/luatex font selection
\fi
% Use upquote if available, for straight quotes in verbatim environments
\IfFileExists{upquote.sty}{\usepackage{upquote}}{}
\IfFileExists{microtype.sty}{% use microtype if available
  \usepackage[]{microtype}
  \UseMicrotypeSet[protrusion]{basicmath} % disable protrusion for tt fonts
}{}
\makeatletter
\@ifundefined{KOMAClassName}{% if non-KOMA class
  \IfFileExists{parskip.sty}{%
    \usepackage{parskip}
  }{% else
    \setlength{\parindent}{0pt}
    \setlength{\parskip}{6pt plus 2pt minus 1pt}}
}{% if KOMA class
  \KOMAoptions{parskip=half}}
\makeatother
\usepackage{xcolor}
\usepackage[margin=1in]{geometry}
\usepackage{longtable,booktabs,array}
\usepackage{calc} % for calculating minipage widths
% Correct order of tables after \paragraph or \subparagraph
\usepackage{etoolbox}
\makeatletter
\patchcmd\longtable{\par}{\if@noskipsec\mbox{}\fi\par}{}{}
\makeatother
% Allow footnotes in longtable head/foot
\IfFileExists{footnotehyper.sty}{\usepackage{footnotehyper}}{\usepackage{footnote}}
\makesavenoteenv{longtable}
\usepackage{graphicx}
\makeatletter
\newsavebox\pandoc@box
\newcommand*\pandocbounded[1]{% scales image to fit in text height/width
  \sbox\pandoc@box{#1}%
  \Gscale@div\@tempa{\textheight}{\dimexpr\ht\pandoc@box+\dp\pandoc@box\relax}%
  \Gscale@div\@tempb{\linewidth}{\wd\pandoc@box}%
  \ifdim\@tempb\p@<\@tempa\p@\let\@tempa\@tempb\fi% select the smaller of both
  \ifdim\@tempa\p@<\p@\scalebox{\@tempa}{\usebox\pandoc@box}%
  \else\usebox{\pandoc@box}%
  \fi%
}
% Set default figure placement to htbp
\def\fps@figure{htbp}
\makeatother
\setlength{\emergencystretch}{3em} % prevent overfull lines
\providecommand{\tightlist}{%
  \setlength{\itemsep}{0pt}\setlength{\parskip}{0pt}}
\setcounter{secnumdepth}{5}
\usepackage[utf8]{inputenc}
\usepackage[bookmarks=false]{hyperref}
\usepackage{xcolor}
\usepackage{pifont}
\usepackage{bookmark}
\IfFileExists{xurl.sty}{\usepackage{xurl}}{} % add URL line breaks if available
\urlstyle{same}
\hypersetup{
  pdftitle={Relatório de Cobertura CNCT--CBO},
  pdfauthor={Equipe BM -- FGV/DGPE},
  hidelinks,
  pdfcreator={LaTeX via pandoc}}

\title{Relatório de Cobertura CNCT--CBO}
\usepackage{etoolbox}
\makeatletter
\providecommand{\subtitle}[1]{% add subtitle to \maketitle
  \apptocmd{\@title}{\par {\large #1 \par}}{}{}
}
\makeatother
\subtitle{Listagem de CBOs que não aparecem na seleção feita pela IA}
\author{Equipe BM -- FGV/DGPE}
\date{18/07/2025}

\begin{document}
\maketitle

\textbf{CBOs não correspondidos e seus perfis (presentes no CNCT, ausentes na predição da IA)}

\begin{longtable}[]{@{}
  >{\raggedright\arraybackslash}p{(\linewidth - 4\tabcolsep) * \real{0.1236}}
  >{\raggedright\arraybackslash}p{(\linewidth - 4\tabcolsep) * \real{0.4382}}
  >{\raggedright\arraybackslash}p{(\linewidth - 4\tabcolsep) * \real{0.4382}}@{}}
\caption{Cursos com CBOs Não Correspondidos}\tabularnewline
\toprule\noalign{}
\begin{minipage}[b]{\linewidth}\raggedright
IDX\_EIXCUR
\end{minipage} & \begin{minipage}[b]{\linewidth}\raggedright
Denominacao\_Curso\_CNCT
\end{minipage} & \begin{minipage}[b]{\linewidth}\raggedright
unmatched\_cbos
\end{minipage} \\
\midrule\noalign{}
\endfirsthead
\toprule\noalign{}
\begin{minipage}[b]{\linewidth}\raggedright
IDX\_EIXCUR
\end{minipage} & \begin{minipage}[b]{\linewidth}\raggedright
Denominacao\_Curso\_CNCT
\end{minipage} & \begin{minipage}[b]{\linewidth}\raggedright
unmatched\_cbos
\end{minipage} \\
\midrule\noalign{}
\endhead
\bottomrule\noalign{}
\endlastfoot
01007 & Técnico em Enfermagem & 322235 \\
01033 & Técnico em Vigilância em Saúde & 352205 \\
02001 & Técnico em Automação Industrial & 313215 \\
02005 & Técnico em Eletrotécnica & 318705 \\
02006 & Técnico em Fabricação Mecânica & 391125 \\
02009 & Técnico em Instrumentação Industrial & 313405, 313415 \\
04002 & Técnico em Comércio & 354205, 354705 \\
04008 & Técnico em Logística & 342110, 391115 \\
04009 & Técnico em Marketing & 354130 \\
05001 & Técnico em Computação Gráfica & 317120 \\
05003 & Técnico em Informática & 317115, 317120 \\
05009 & Técnico em Telecomunicações & 391205 \\
06005 & Técnico em Edificações & 318005, 318010, 318015, 318120 \\
06011 & Técnico em Portos & 354305 \\
06017 & Técnico em Transporte Metroferroviário & 342120, 342315 \\
06018 & Técnico em Transporte Rodoviário & 342210 \\
09001 & Técnico em Artes Circenses & 376305, 376310, 376315, 376320 \\
09011 & Técnico em Design de Interiores & 318005, 318010 \\
09014 & Técnico em Design de Móveis & 771120 \\
09019 & Técnico em Figurino Cênico & 318810 \\
09023 & Técnico em Multimídia & 317110 \\
09030 & Técnico em Produção de Áudio e Vídeo & 374305, 517330 \\
09031 & Técnico em Produção de Moda & 375110 \\
09033 & Técnico em Rádio e Televisão & 374110, 374115, 374120, 374135,
374140 \\
10001 & Técnico em Açúcar e Álcool & 311115, 325305, 325310 \\
10006 & Técnico em Celulose e Papel & 311115, 311710, 311715 \\
10008 & Técnico em Construção Naval & 318205, 318210 \\
10020 & Técnico em Têxtil & 311705, 311715, 311725, 721305 \\
11016 & Técnico em Recursos Pesqueiros & 321310, 341220 \\
12001 & Técnico em Defesa Civil & 351605 \\
\end{longtable}

\newpage
\section{ 01007 -- Técnico em Enfermagem }

\textbf{Perfil do Curso (CNCT)}

O Técnico em Enfermagem será habilitado para:

Realizar, sob a supervisão do Enfermeiro, cuidados integrais de
enfermagem a indivíduos, família e grupos sociais vulneráveis ou não.
Atuar na promoção, prevenção, recuperação e reabilitação dos processos
saúde-doença em todo o ciclo vital. Participar do planejamento e
execução das ações de saúde junto à equipe multidisciplinar,
considerando as normas de biossegurança, envolvendo curativos,
administração de medicamentos e vacinas, nebulizações, banho de leito,
cuidados pós-morte, mensuração antropométrica e verificação de sinais
vitais. Preparar o paciente para os procedimentos de saúde. Participar
de comissões de certificação de serviços de saúde, tais como núcleo de
segurança do paciente, serviço de controle de infecção hospitalar,
gestão da qualidade dos serviços prestados à população, gestão de
riscos, comissões de ética de enfermagem, transplantes, óbitos e outros.
Colaborar com o Enfermeiro em ações de comissões de certificação de
serviços de saúde, tais como núcleo de segurança do paciente, serviço de
controle de infecção hospitalar, gestão da qualidade dos serviços
prestados à população, gestão de riscos, comissões de ética de
enfermagem, transplantes, óbitos e outros.

Para a atuação como Técnico em Enfermagem, são fundamentais:

Conhecimentos das políticas públicas de saúde e compreensão de sua
atuação profissional frente às diretrizes, princípios e estrutura
organizacional do Sistema Único de Saúde (SUS). Conhecimentos e saberes
relacionados aos princípios das técnicas aplicadas na área, sempre
pautados numa postura humana e ética. Resolução de situações-problema,
comunicação, trabalho em equipe e interdisciplinar, tecnologias da
informação e da comunicação, gestão de conflitos e ética profissional.
Organização e responsabilidade. Iniciativa social. Determinação e
criatividade, buscando promover a humanização da assistência.
Atualização e aperfeiçoamento profissional por meio da educação
continuada.

\textbf{CBOs não correspondidos e seus perfis:}

\textbf{CBO: 322235  ( Auxiliar de enfermagem do trabalho )}

Participa da organização e funcionamento do serviço de segurança e saúde
do trabalhador. Atua em situações de urgência e emergência, integrando a
equipe multidisciplinar de saúde e do serviço especializado em
engenharia e medicina do trabalho (SESMT). Atende, sob supervisão,
trabalhadores portadores de doenças ou lesões de pouca gravidade.
Auxilia na observação sistemática do estado de saúde dos trabalhadores,
da ocorrência de doenças profissionais, lesões traumáticas e doenças
epidemiológicas. Preenche relatório dos dados observados. Auxilia na
realização de exame pré-admissional, periódicos e demissionais. Realiza
ações de prevenção de doenças ocupacionais e promoção da saúde e
segurança no trabalho. Participa de programas de prevenção de acidentes
e de reabilitação profissional. Participa de campanhas de educação
sanitária e de vacinação. Pode atuar em programas de qualidade e de
certificação de serviços de saúde do trabalhador.

\begin{Form}
\textbf{Esta correspondência está adequada?} \\ 
\ChoiceMenu[radio,radiosymbol=\ding{52},name=cbovalid_ 322235 ]{}{CBO deve ficar=sim, \hspace{6em} CBO não fica aqui=nao}
\end{Form}

\newpage
\section{ 01033 -- Técnico em Vigilância em Saúde }

\textbf{Perfil do Curso (CNCT)}

O Técnico em Vigilância em Saúde será habilitado para:

Atuar, sob a supervisão de profissionais de nível superior, de forma a
aplicar a normatização de produtos, processos, ambientes e serviços de
interesses a saúde. Controlar o fluxo de pessoas, animais, plantas e
produtos em portos, aeroportos e fronteiras. Desenvolver ações de
controle e monitoramento de doenças, endemias e de vetores. Desenvolver
ações de inspeção e fiscalização sanitárias. Elaborar e implementar,
junto com a população do território, ações educativas no âmbito das
vigilâncias: ambiental, sanitária, epidemiológica e saúde do trabalhador
para promoção da saúde. Investigar, monitorar e avaliar riscos e os
determinantes dos agravos e danos à saúde da população, relacionados ao
trabalho, assim como ao meio ambiente. Planejar e articular ações
intersetoriais para promoção da saúde prevenção, controle e
monitoramento de vetores e doenças endêmicas. Planejar, executar e
avaliar o processo de vigilância sanitária, epidemiológica, ambiental,
saúde do trabalhador e vigilância da situação de saúde. Realizar análise
territorial das condições de vida e de saúde da população, como também
identificar e intervir em situações de risco e de vulnerabilidade de
grupos populacionais e ambientes. Desenvolver ações de mobilização
social junto à comunidade para promoção e proteção da saúde. Realizar
ações de prevenção e controle de doenças e agravos à saúde em interação
com o Agente Comunitário de Saúde e a equipe de Atenção Básica.

Para a atuação como Técnico em Vigilância em Saúde, são fundamentais:

Conhecimento das políticas públicas de saúde: organização, princípios e
diretrizes do Sistema Único de Saúde (SUS). Conhecimentos e saberes
relacionados a processos de sustentabilidade, territorialização,
organização do SUS e vigilâncias. Resolução de situações-problemas,
trabalho em equipe, tecnologias da informação e da comunicação, gestão
de conflitos e ética profissional. Iniciativa e capacidade de
planejamento Organização, responsabilidade, iniciativa social,
entusiasmo, empatia e respeito. Atualização e aperfeiçoamento
profissional por meio da educação continuada.

\textbf{CBOs não correspondidos e seus perfis:}

\textbf{CBO: 352205  ( Agente de defesa ambiental )}

Protege a fauna e a flora, impedindo as práticas que provoquem a
extinção de espécies ou submetam os animais à crueldade. Fiscaliza
atividades e obras, solicitando documentação ao fiscalizado e coletando
informações técnicas relacionadas à preservação do meio ambiente.
Vistoria atividades e obras, verificando o cumprimento das exigências
legais e técnicas e avaliando o impacto ambiental. Analisa tecnicamente
projetos, examinando o processo de licenciamento, enviando material para
análise nos órgãos competentes e elaborando parecer, laudo ou relatório
técnico. Investiga denúncias de danos ao meio ambiente, levantando
informações junto à comunidade local da ocorrência e coletando material
para análise. Se constatada a veracidade da denúncia, enquadra
legalmente o caso. Reprime atividades ilegais, interditando
estabelecimentos e atividades, embargando obras, apreendendo animais em
condições ilegais e providenciando destinação de produtos irregulares.
Adverte, notifica, intima ou multa infrator, podendo encaminhá-lo às
autoridades competentes. Toma providências para minimizar impactos de
acidentes ambientais, de destruições e de outros danos causados em
recursos naturais e na biodiversidade. Previne e combate incêndios.
Presta suporte técnico à polícia florestal. Ministra palestras e promove
educação ambiental em comunidades, escolas e espaços públicos, para
sensibilizar crianças, jovens e adultos para a importância da defesa do
meio ambiente. Pode gerenciar atividades, coordenando e supervisionando
equipes e administrando recursos humanos, financeiros e materiais.
Providencia manutenção de equipamentos e instalações. Emite
notificações, autorizações, intimações, licenças e ofícios. Registra
denúncias, solicita mandado de busca e apreensão e formaliza propostas
de embargo, interdição e multa. Consulta imagens e relatórios de
geoinformação para realizar o seu trabalho.

\begin{Form}
\textbf{Esta correspondência está adequada?} \\ 
\ChoiceMenu[radio,radiosymbol=\ding{52},name=cbovalid_ 352205 ]{}{CBO deve ficar=sim, \hspace{6em} CBO não fica aqui=nao}
\end{Form}

\newpage
\section{ 02001 -- Técnico em Automação Industrial }

\textbf{Perfil do Curso (CNCT)}

O Técnico em Automação Industrial será habilitado para:

Desenvolver e integrar soluções para sistemas de automação visando à
medição e ao controle de variáveis em processos industriais,
considerando as normas, os padrões e os requisitos técnicos de
qualidade, saúde e segurança e de meio ambiente. Empregar programas de
computação e redes industriais no controle de processos industriais.
Planejar, controlar e executar a instalação e a manutenção de
equipamentos automatizados e/ou sistemas robotizados para controle de
processos industriais. Realizar medições, testes e calibrações em
equipamentos eletroeletrônicos empregados em controle de processos
industriais. Instalar, configurar e operar tecnologias de manufatura
aditiva, sistemas ciberfísicos e processos de produção com internet das
coisas. Reconhecer tecnologias inovadoras presentes no segmento visando
a atender às transformações digitais na sociedade. Realizar
especificação, projeto, instalação, medição, teste, diagnóstico e
calibração de equipamentos e sistemas automatizados. Executar
procedimentos de controle de qualidade, operação e gestão de sistemas
automatizados e controle de processos.

Para atuação como Técnico em Automação Industrial, são fundamentais:

Conhecimentos e saberes relacionados aos processos de planejamento e
implementação de processos automatizados de modo a assegurar a saúde e a
segurança dos trabalhadores e dos usuários. Conhecimentos e saberes
relacionados à sustentabilidade do processo produtivo, às técnicas e aos
processos de produção, às normas técnicas, à liderança de equipes, à
solução de problemas técnicos e trabalhistas e à gestão de conflitos.

\textbf{CBOs não correspondidos e seus perfis:}

\textbf{CBO: 313215  ( Técnico eletrônico )}

Conserta e faz manutenção corretiva em aparelhos e equipamentos
eletrônicos, interpretando esquemas elétricos e substituindo componentes
defeituosos. Pode utilizar equipamentos de diagnóstico eletrônico.
Instala equipamentos, aparelhos e dispositivos eletrônicos, observando
as condições do ambiente, realizando ajustes e testes. Faz manutenções
preventivas e preditivas em equipamentos eletrônicos, de acordo com
planos de manutenção, podendo fazer uso da internet para suporte.
Realiza medições, testes e calibrações de equipamentos. Desenvolve,
monta e testa dispositivos de circuitos eletrônicos, especificando
componentes, calculando custos, incluindo a possibilidade de uso de
microcontroladores, microprocessadores e outros dispositivos eletrônicos
digitais programáveis. Sugere alterações no processo de produção,
implementando dispositivos de automação e integração de dispositivos e
sistemas eletrônicos por meio de redes digitais de comunicação de dados,
locais e a distância. Participa de reuniões técnicas com pessoal interno
e externo. Redige procedimentos de trabalho e elabora gráficos de
resultados. Pode realizar treinamento, orientação e avaliação de
desempenho de operadores.

\begin{Form}
\textbf{Esta correspondência está adequada?} \\ 
\ChoiceMenu[radio,radiosymbol=\ding{52},name=cbovalid_ 313215 ]{}{CBO deve ficar=sim, \hspace{6em} CBO não fica aqui=nao}
\end{Form}

\newpage
\section{ 02005 -- Técnico em Eletrotécnica }

\textbf{Perfil do Curso (CNCT)}

O Técnico em Eletrotécnica será habilitado para:

Planejar, controlar e executar a instalação e a manutenção de sistemas e
instalações elétricas industriais, prediais e residenciais, considerando
as normas, os padrões e os requisitos técnicos de qualidade, saúde e
segurança e de meio ambiente. Elaborar e desenvolver projetos de
instalações elétricas industriais, prediais e residenciais, sistemas de
acionamentos elétricos e de automação industrial e de infraestrutura
para sistemas de telecomunicações em edificações. Aplicar medidas para o
uso eficiente da energia elétrica e de fontes energéticas alternativas.
Elaborar e desenvolver programação e parametrização de sistemas de
acionamentos eletrônicos industriais. Planejar e executar instalação e
manutenção de sistemas de aterramento e de descargas atmosféricas em
edificações residenciais, comerciais e industriais. Reconhecer
tecnologias inovadoras presentes no segmento visando a atender às
transformações digitais na sociedade.

Para atuação como Técnico em Eletrotécnica, são fundamentais:

Conhecimentos e saberes relacionados aos processos de planejamento e
implementação de sistemas elétricos de modo a assegurar a saúde e a
segurança dos trabalhadores e dos usuários. Conhecimentos e saberes
relacionados à sustentabilidade do processo produtivo, às técnicas e aos
processos de produção, às normas técnicas, à liderança de equipes, à
solução de problemas técnicos e trabalhistas e à gestão de conflitos.

\textbf{CBOs não correspondidos e seus perfis:}

\textbf{CBO: 318705  ( Desenhista projetista de eletricidade )}

Participa da elaboração de anteprojetos elétricos, efetuando estudos,
elaborando desenhos para apreciação do cliente e definindo os pontos
críticos do projeto. Define os recursos humanos, estima sua quantidade e
participa da distribuição e da atribuição de responsabilidades. Realiza
as estimativas de materiais. Participa da elaboração do cronograma de
desenvolvimento e estabelece o prazo de entrega. Acompanha o
desenvolvimento do projeto. Analisa especificações técnicas fornecidas
pelo cliente e coleta dados -- utilizando técnicas de pesquisa e fontes
diversificadas -- para elaboração do projeto. Examina normas técnicas
relacionadas ao projeto. Desenvolve projetos elétricos de produtos e
instalações elétricas, utilizando ferramentas computacionais integradas
na elaboração dos desenhos e demais documentos técnicos do projeto.
Realiza cálculos técnicos de projetos elétricos; desenha diagramas
elétricos detalhados; e simula o projeto elétrico, por meio de programas
aplicativos específicos para eletricidade do tipo CAD-Desenho Assistido
por Computador (Computer Aided Design). Pode utilizar programas
aplicativos do tipo CAM--Manufatura Auxiliada por Computador (Computer
Aided Manufacturing), para geração de programas e informações de
prototipagem ou de fabricação -- de placas de circuito impresso, por
exemplo --, quando requeridas no projeto. Utiliza normas de qualidade e
de segurança do produto, nas especificações e nos cálculos do projeto.
Faz especificação de componentes, módulos e outros materiais elétricos
aplicáveis ao projeto. Incorpora componentes eletrônicos ao projeto - a
fim de agregar valor e novas funções ou reduzir custos - e introduz
inovações relacionadas à eficiência energética. Otimiza os custos
envolvidos em todas as partes do projeto. Adapta o projeto em função das
necessidades do cliente. Confere o projeto e faz correções ou
modificações, quando necessário. Elabora lista de materiais, folha de
dados (datasheet), manuais de instalação e de manutenção, catálogos e
outros documentos técnicos do projeto. Participa de reuniões para tomada
de decisão e coordena equipes de trabalho na elaboração do projeto. Zela
pelo cumprimento do cronograma do projeto. Desenvolve fornecedores para
os projetos elétricos de produtos e instalações elétricas. Pesquisa
fornecedores em potencial, avaliando-os por meio de visitas técnicas.
Elabora desenhos para produtos especiais -- tais como componentes,
peças, módulos, equipamentos ou acessórios para projetos elétricos --
e/ou introduz melhorias em produtos especiais, junto aos fornecedores.
Acompanha etapas de desenvolvimento de projetos e ensaios de amostras.
Avalia a relação custo-benefício dos produtos fornecidos. Participa da
implantação do projeto elétrico de produtos e instalações elétricas,
auxiliando no desenvolvimento do protótipo; colaborando na definição de
métodos de produção; e acompanhando fabricação e montagem. Introduz
modificações no projeto, em função de erros detectados e necessidades de
implementação de melhorias. Participa de reuniões de avaliação do
projeto. Acompanha ensaios finais de produtos e instalações elétricas
projetadas, estabelecendo tipos de equipamentos e parâmetros para
testes; analisando resultados de testes; e elaborando relatórios de
ensaios. Modifica o projeto, em função dos resultados de testes.
Controla documentação relacionada ao projeto, estruturando e mantendo
arquivos -- físico e eletrônico --; atualizando cópias do projeto; e
registrando etapas e revisões de projeto. Nos processos de
desenvolvimento de projetos elétricos de produtos e instalações
elétricas, aplica conceitos e ferramentas da qualidade. Conserva a
limpeza, controla a luminosidade e verifica as condições ergonômicas do
ambiente de trabalho. Coleta e descarta resíduos, de acordo com
diretrizes de preservação ambiental. Conserva equipamentos em plenas
condições de funcionamento. Segue preceitos de segurança no trabalho,
utilizando equipamentos de proteção individual.

\begin{Form}
\textbf{Esta correspondência está adequada?} \\ 
\ChoiceMenu[radio,radiosymbol=\ding{52},name=cbovalid_ 318705 ]{}{CBO deve ficar=sim, \hspace{6em} CBO não fica aqui=nao}
\end{Form}

\newpage
\section{ 02006 -- Técnico em Fabricação Mecânica }

\textbf{Perfil do Curso (CNCT)}

O Técnico em Fabricação Mecânica será habilitado para:

Desenvolver projetos, planejar, supervisionar e controlar atividades de
fundição, em usinagem convencional e computadorizada, em caldeiraria, em
soldagem e processos de conformação mecânica. Interpretar desenho
técnico. Selecionar, desenvolver e especi?car ferramental para os
processos produtivos. Executar ensaios mecânicos. Especi?car materiais e
insumos aplicados aos processos de fabricação mecânica. Controlar
estoques de produtos acabados.

Para atuação como Técnico em Fabricação Mecânica, são fundamentais:

Conhecimentos e saberes relacionados aos processos de planejamento e
operação das atribuições da área de modo a assegurar a saúde e a
segurança dos trabalhadores e dos futuros usuários e operadores de
empresas em processos de transformação em fabricação mecânica.
Conhecimentos e saberes relacionados à sustentabilidade do processo
produtivo, às normas e relatórios técnicos, à legislação da área, às
novas tecnologias relacionadas à indústria 4.0, à liderança de equipes,
à solução de problemas técnicos e à gestão de conflitos.

\textbf{CBOs não correspondidos e seus perfis:}

\textbf{CBO: 391125  ( Técnico de planejamento de produção )}

Planeja a produção, a partir da análise da demanda e da capacidade de
produção. Realiza a programação da produção, analisando ordens de
produção, determinando prioridades, definindo cronograma e estabelecendo
a sequência de produção. Estabelece parâmetros para controle dos
processos de produção. Controla o desenvolvimento dos planos de
produção, com base em parâmetros de controle e monitoramento do fluxo
produtivo. Identifica desvios e define ações corretivas. Propõe
melhorias no processo de produção. Define estoque de segurança de
suprimentos necessários para atender às demandas de produção.
Providencia entrega e distribuição de suprimentos para agilizar o fluxo
de materiais e cumprir os cronogramas de produção. Pode revisar os
cronogramas de produção por falta de trabalhadores, atraso de entrega de
material ou outros problemas. Trata as informações do ciclo
planejamento-produção, elaborando gráficos e relatórios de controle.
Disponibiliza e atualiza informações. Interage com profissionais
envolvidos nas diversas etapas de produção, com clientes e com
fornecedores, para coordenar as atividades de produção, para resolver
reclamações ou para eliminar atrasos. Auxilia na definição das
instruções de trabalho.

\begin{Form}
\textbf{Esta correspondência está adequada?} \\ 
\ChoiceMenu[radio,radiosymbol=\ding{52},name=cbovalid_ 391125 ]{}{CBO deve ficar=sim, \hspace{6em} CBO não fica aqui=nao}
\end{Form}

\newpage
\section{ 02009 -- Técnico em Instrumentação Industrial }

\textbf{Perfil do Curso (CNCT)}

O Técnico em Instrumentação Industrial será habilitado para:

Implantar, configurar e manter sistemas de instrumentação e controle de
processos industriais. Elaborar, desenvolver e executar projetos de
instalações de instrumentos de medição de variáveis industriais, de
acordo com normas técnicas, ambientais, de saúde e segurança no trabalho
e de qualidade. Realizar programação, parametrização, medições e testes
de equipamentos por meio de instrumentação industrial. Reconhecer
tecnologias inovadoras presentes no segmento visando a atender às
transformações digitais na sociedade.

Para atuação como Técnico em Instrumentação Industrial, são
fundamentais:

Conhecimentos e saberes relacionados aos processos de planejamento e
implementação de controle de processos industriais por meio de
instrumentação de modo a assegurar a saúde e a segurança dos
trabalhadores. Conhecimentos e saberes relacionados à sustentabilidade
do processo produtivo, às técnicas e processos de produção, às normas
técnicas, à liderança de equipes, à solução de problemas técnicos e
trabalhistas e à gestão de conflitos

\textbf{CBOs não correspondidos e seus perfis:}

\textbf{CBO: 313405  ( Técnico em calibração )}

Prepara o serviço de calibração, identificando especificações técnicas,
monitorando as condições ambientais para a calibração, inspecionando
visualmente os itens a serem calibrados e determinando o procedimento de
calibração. Realiza aferição e calibração de padrões, equipamentos,
sistemas e instrumentos de medição e controle. Calcula e interpreta os
resultados das medições realizadas no processo, comparando com padrões,
para realizar os ajustes necessários nos equipamentos, sistemas e
instrumentos. Realiza medições de grandezas, lendo e interpretando
desenhos técnicos, selecionando métodos e princípios de medição. Executa
ensaios físicos, elétricos, mecânicos e eletrônicos, dentre outros,
correspondentes às grandezas a serem medidas. Opera padrões,
equipamentos, sistemas e instrumentos de medição e controle. Emite
laudos e certificados da calibração de equipamentos e instrumentos de
medição e controle. Arquiva documentação técnica e gerencial. Participa
da gestão dos procedimentos do sistema de confiabilidade metrológica de
laboratórios de calibração credenciados, elaborando sistema de
codificação, estabelecendo frequência, controlando prazos e validando
resultados da calibração. Participa de auditorias internas e externas.
Participa da gestão da documentação técnica de laboratórios
credenciados, elaborando, atualizando e analisando criticamente
procedimentos e instruções de trabalho, cadastrando instrumentos de
medição e controle, elaborando fichas e formulários de registros.
Contribui no desenvolvimento de projetos de sistemas de medição e
controle, identificando as variáveis envolvidas no processo e avaliando
o desempenho de instrumentos e sistemas em desenvolvimento, realizando a
calibração de protótipos e novos produtos. Realiza manutenção de
instrumentos de medição e controle, em bancada, para fins de calibração,
identificando disfunções e suas causas, relacionando custos e benefícios
da manutenção, reparando ou substituindo componentes danificados.
Realiza a calibração dos instrumentos após os trabalhos de manutenção.

\begin{Form}
\textbf{Esta correspondência está adequada?} \\ 
\ChoiceMenu[radio,radiosymbol=\ding{52},name=cbovalid_ 313405 ]{}{CBO deve ficar=sim, \hspace{6em} CBO não fica aqui=nao}
\end{Form}

\textbf{CBO: 313415  ( Encarregado de manutenção de instrumentos de controle, medição e similares )}

Planeja, com os integrantes da equipe de trabalho, a execução de
manutenção preventiva, preditiva e corretiva de instrumentos elétricos,
eletrônicos e mecânicos de medição, controle e funções análogas,
utilizados em processos industriais. Supervisiona as atividades de
manutenção, estabelecendo procedimentos para identificar disfunções em
instrumentos, por meio de sistemas manuais, semiautomáticos e
automáticos. Contribui no aperfeiçoamento de sistemas de medição e
controle, analisando indicadores das atividades de manutenção de
instrumentos -- tais como TMEF-Tempo Médio Entre Falhas (MTBF--Mean Time
Between Failures), TMPR-Tempo Médio Para Reparo (MTTR-Mean Time To
Repair), dentre outros -, para calcular a taxa de disponibilidade dos
equipamentos e do processo; retroalimentar o sistema; avaliar seu
desempenho; e propor soluções aperfeiçoadas de medição e controle, com
base em dados e informações. Garante a calibração dos instrumentos,
monitorando a validade de calibração e selecionando os prestadores de
serviços. Pode supervisionar as atividades de calibração em laboratórios
credenciados. Participa da gestão dos procedimentos do sistema de
confiabilidade metrológica, colaborando na elaboração de sistema de
codificação, estabelecendo frequência, controlando prazos e validando os
resultados da calibração. Participa de auditorias internas e externas.
Analisa tecnicamente a aquisição de produtos e serviços para manutenção,
estabelecendo os objetivos da análise, indicando fornecedores potenciais
e definindo características técnicas de produtos e serviços. Emite
pareceres técnicos. Participa da gestão de contratos de serviços
terceirizados de manutenção. Organiza a documentação técnica e
gerencial, elaborando, atualizando e analisando criticamente
procedimentos e instruções de trabalho, cadastrando instrumentos e
elaborando fichas e formulários de registros. Coordena e supervisiona a
equipe de trabalho, orientando os trabalhos, avaliando os desempenhos
individuais e coletivos e identificando as necessidades de treinamento.

\begin{Form}
\textbf{Esta correspondência está adequada?} \\ 
\ChoiceMenu[radio,radiosymbol=\ding{52},name=cbovalid_ 313415 ]{}{CBO deve ficar=sim, \hspace{6em} CBO não fica aqui=nao}
\end{Form}

\newpage
\section{ 04002 -- Técnico em Comércio }

\textbf{Perfil do Curso (CNCT)}

O Técnico em Comércio será habilitado para:

Aplicar métodos de comercialização de bens e serviços em loja física ou
virtual. Efetuar controle quantitativo e qualitativo de produtos, preços
e tributos. Coordenar e controlar a armazenagem em estabelecimento
comercial. Elaborar planilha de custos. Identificar demanda e comunicar
previsões a fornecedores. Ofertar serviços correlatos aos produtos
comercializados. Operacionalizar planos de marketing e de comunicação.
Executar atividades voltadas à logística, a recursos humanos e à
comercialização.

Para atuação como Técnico em Comércio, são fundamentais:

Conhecimentos e saberes relacionados ao funcionamento da área comercial
e de prestação de serviços, de modo a atuar em conformidade com as
legislações e diretrizes de conduta, como também com as normas de saúde
e segurança do trabalho. Atuação de forma proativa em atividades de
comercialização de produtos e serviços, com visão empreendedora,
comunicação clara e cordial, comprometimento com necessidades e desejos
de clientes e respeito a demais stakeholders.

\textbf{CBOs não correspondidos e seus perfis:}

\textbf{CBO: 354205  ( Comprador )}

Ao receber requisições de compras, verifica as especificações dos
materiais ou dos serviços. Pode agrupar os itens requeridos por lotes
(famílias ou ramos). Verifica quantidade de materiais solicitados,
informando-se sobre a disponibilidade em estoques. Consulta a existência
de verba disponível para a aquisição de materiais ou serviços. Realiza a
compra por processo de cotação ou por processo licitatório, levando em
conta - para optar por um deles - o valor estimado da aquisição e a
natureza da empresa ou instituição (pública ou privada). Executa
processo de cotação, determinando - aos fornecedores interessados - a
data para apresentação de propostas. Monta e analisa planilha de
cotações. Pode negociar, com fornecedores, preços, prazos de entrega e
condições de pagamentos. Realiza ou providencia a realização de processo
licitatório, estabelecendo em edital qual será a modalidade e os
requisitos para que os interessados possam inscrever-se. Ao final,
solicita, à área jurídica, a elaboração e a celebração de contrato com o
fornecedor. Concretiza a compra de materiais ou serviços, após
autorização. Acompanha o fluxo de entrega, conferindo se ocorre dentro
das condições estabelecidas. Pode identificar defeitos ou problemas,
rejeitando e devolvendo materiais. Pode negociar substituição ou
reparação dos prejuízos decorrentes da devolução dos materiais.
Desenvolve fornecedores de materiais e serviços, entrevistando,
pesquisando referências, requisitando amostras ou catálogos de materiais
ou serviços e agendando visita técnica ao fornecedor. Pode participar da
homologação de fornecedores. Cadastra e atualiza constantemente dados e
informações de fornecedores. Prepara relatórios de visitas a
fornecedores, de não conformidade de materiais e de variação de preços.
Redige comunicado aos setores sobre compras efetuadas e elabora
correspondências. Organiza e arquiva documentos dos processos de
compras. Pode utilizar sistemas, aplicativos e plataformas, além do
sistema informatizado de controle de compras. Entre outros, pode
utilizar o Sistema Integrado de Gestão Empresarial (Enterprise Resource
Planning -- ERP) - para obter informações de estoque e de existência de
verba disponível -, a plataforma de Gestão de Relacionamento com o
Cliente (Customer Relationship Management - CRM) - para inserir e
consultar dados e informações cadastrais de fornecedores - e o
Marketplace virtual - para verificar preços e prazo de entrega de
diferentes fornecedores -.

\begin{Form}
\textbf{Esta correspondência está adequada?} \\ 
\ChoiceMenu[radio,radiosymbol=\ding{52},name=cbovalid_ 354205 ]{}{CBO deve ficar=sim, \hspace{6em} CBO não fica aqui=nao}
\end{Form}

\textbf{CBO: 354705  ( Representante comercial autônomo )}

Realiza planejamento das atividades de representação comercial autônoma
- de empresas nacionais ou estrangeiras -, fazendo pesquisa de mercado,
analisando potencial de clientes e elaborando plano estratégico.
Estabelece metas de venda de produtos e serviços. Realiza divulgação de
produtos e serviços, definindo estratégias de marketing e requisitando
material de divulgação. Envia mala direta ao conjunto de clientes e
divulga produtos por meio de multimídia. Informa aos clientes sobre
alterações nos produtos e serviços, ressaltando as melhorias. Prepara
cadastro de novos clientes e estabelece contato com eles por telefone,
celular, e-mail e outros meios de comunicação. Realiza vendas aos novos
clientes. Planeja itinerário de visitas, agendando horários. Explica os
objetivos da visita, demonstra benefícios e qualidades de produtos e
serviços e fornece amostras. Auxilia na escolha, prestando informações
técnicas dos produtos. Apresenta tabelas de preços, explica formas de
pagamento e efetiva vendas. Recebe clientes para vendas, identificando
suas necessidades e apresentando proposta ao cliente. Calcula custo do
produto e serviço. Estabelece prazos de entrega dos produtos. Finaliza
as vendas, emitindo pedidos de produtos e serviços. Esclarece dúvidas
dos clientes referentes aos pedidos e aos contratos de prestação de
serviços. Monitora o cumprimento do prazo de entrega dos produtos.
Acompanha clientes pós-venda, verificando suas opiniões. Filtra
informações para melhoria de serviços, dando retorno às sugestões dos
clientes e enviando-lhes brindes. Renova e altera contratos de prestação
de serviços. Interage com as diversas áreas da empresa representada.
Requisita mostruário de produtos e solicita troca de produtos
defeituosos. Obtém informações de pré-venda com telemarketing. Comunica
transações comerciais realizadas e os resultados de vendas. Sugere
políticas de vendas, propondo ações para melhoria dos resultados. Visita
área industrial da empresa, participando de reuniões sobre promoção de
produtos e serviços. Mantém-se atualizado, participando de cursos sobre
novas estratégias de marketing. Participa de eventos, tais como feiras e
palestras relacionadas à área. Expõe produtos em locais definidos pela
empresa representada.

\begin{Form}
\textbf{Esta correspondência está adequada?} \\ 
\ChoiceMenu[radio,radiosymbol=\ding{52},name=cbovalid_ 354705 ]{}{CBO deve ficar=sim, \hspace{6em} CBO não fica aqui=nao}
\end{Form}

\newpage
\section{ 04008 -- Técnico em Logística }

\textbf{Perfil do Curso (CNCT)}

O Técnico em Logística será habilitado para:

Auxiliar no planejamento, operacionalização e controle da cadeia
produtiva e seu fluxo logístico. Executar procedimentos relacionados a
suprimentos, produção, recebimento, armazenagem e distribuição de
produtos, fazendo uso das tecnologias de informação e comunicação.
Identificar agentes da cadeia de suprimentos. Elaborar relatórios
operacionais para tomada de decisões.

Para atuação como Técnico em Logística, são fundamentais:

Conhecimentos e saberes relacionados à área operacional, de produção e
de prestação de serviços das organizações, de modo a atuar em
conformidade com as legislações e diretrizes de conduta, como também com
as normas de saúde e segurança do trabalho. Atuação de forma proativa na
resolução de situações-problema do mundo do trabalho, desenvolvendo
competências socioemocionais e atributos comportamentais relacionados à
sustentabilidade e ao trabalho colaborativo.

\textbf{CBOs não correspondidos e seus perfis:}

\textbf{CBO: 342110  ( Operador de transporte multimodal )}

Programa, controla e coordena operações de transportes em geral - em
diferentes modais de transporte -, acompanhando as operações de
embarque, transbordo e desembarque de carga junto ao regime de trânsito
aduaneiro. Supervisiona as condições de segurança e de manutenção dos
meios de transportes e equipamentos utilizados. Controla a eficiência
operacional do armazenamento e do transporte da carga. Remaneja
equipamentos e veículos, quando necessário. Controla recursos
financeiros e insumos, elaborando planilhas de custos e receitas dos
serviços prestados. Pesquisa mercado, analisando a viabilidade de novos
serviços, coletando dados de clientes e fornecedores potenciais e
identificando rotas de transporte e respectiva estrutura de serviço.
Define parcerias e fornecedores de serviços com critérios técnicos,
contratando serviços de transporte e negociando preços e prazos de
pagamento. Elabora ou auxilia na elaboração de documentação - como
contratos, registros legais de transporte de cargas do comércio
doméstico ou do exterior - necessária ao desembaraço de cargas. Registra
dados em sistemas e alimenta sistemas informacionais internos. Atende
clientes, classificando-os por demanda, identificando suas necessidades,
elaborando propostas de preços e serviços, disponibilizando informações
da carga e verificando a sua satisfação. Treina e orienta funcionários e
estagiários, avaliando os seus desempenhos e analisando as necessidades
de treinamento da equipe.

\begin{Form}
\textbf{Esta correspondência está adequada?} \\ 
\ChoiceMenu[radio,radiosymbol=\ding{52},name=cbovalid_ 342110 ]{}{CBO deve ficar=sim, \hspace{6em} CBO não fica aqui=nao}
\end{Form}

\textbf{CBO: 391115  ( Controlador de entrada e saída )}

Controla suprimentos - matéria-prima e insumos - nos processos de
produção, registrando entrada e saída, definindo estoque de segurança
para cada item e emitindo solicitações. Pode propor suprimentos
alternativos. Define formas de transporte, manuseio, armazenamento e
distribuição de suprimentos. Realiza a inspeção de matéria-prima e
insumos. Mantém registros relacionados ao movimento de entrada e saída
de suprimentos. Atualiza, disponibiliza e trata as informações,
elaborando gráficos e relatórios de controles. Colabora no planejamento
e no controle da produção, propondo aperfeiçoamentos nos processos,
indicando melhorias em máquinas e equipamentos e definindo leiaute do
processo produtivo. Pode auxiliar na definição de instruções de
trabalho. Mantém comunicação com clientes e fornecedores.

\begin{Form}
\textbf{Esta correspondência está adequada?} \\ 
\ChoiceMenu[radio,radiosymbol=\ding{52},name=cbovalid_ 391115 ]{}{CBO deve ficar=sim, \hspace{6em} CBO não fica aqui=nao}
\end{Form}

\newpage
\section{ 04009 -- Técnico em Marketing }

\textbf{Perfil do Curso (CNCT)}

O Técnico em Marketing será habilitado para:

Projetar e implementar planos de marketing. Realizar análises de vendas,
preços e produtos. Desenvolver projetos de comunicação, ?delização de
clientes e relação com fornecedores ou outras entidades. Desenvolver,
implementar e gerenciar estratégias de marketing digital.
Operacionalizar apresentação dos serviços e produtos no ponto de venda.
Elaborar estudos de mercado.

Para atuação como Técnico em Marketing, são fundamentais:

Conhecimentos e saberes relacionados à área comercial e de negócios das
organizações, de modo a atuar em conformidade com as legislações e
diretrizes de conduta, como também com as normas de saúde e segurança do
trabalho, demonstrando visão empreendedora. Competências socioemocionais
e atributos comportamentais relacionados à comunicação clara e cordial,
respeito à diversidade, trabalho colaborativo e protagonismo na análise
e resolução de problemas voltados ao mundo do trabalho.

\textbf{CBOs não correspondidos e seus perfis:}

\textbf{CBO: 354130  ( Promotor de vendas especializado )}

Promove o processo de venda de produtos especializados - inclusive por
comércio eletrônico (ou e-commerce) -, visando ao aumento de vendas e à
fixação de marca. Propõe política de vendas para produtos
especializados, definindo estratégia de marketing e fazendo previsão de
vendas. Subsidia a revisão de plano de vendas, sugerindo melhorias no
processo promocional e comercial e realizando sondagens de potencial de
mercado. Pesquisa compradores e revendedores potenciais, mapeando-os por
georreferenciamento. Pode elaborar roteiro de visitas, utilizando banco
de dados de compradores e revendedores. Organiza documentação e elabora
manuais técnicos - em meio impresso e digital - para a promoção de
vendas junto a revendedores e grandes compradores, conforme
procedimentos de marketing e vendas. Divulga e demonstra produtos
especializados, de acordo com estratégias de marketing e especificidades
de cada linha desses produtos. Presta assistência a revendedores e a
grandes compradores, instruindo-os sobre o manuseio ou a aplicação dos
produtos. Soluciona problemas para usos especiais, por meio de consulta
ao setor especializado de projeto ou de fabricação. Recebe, registra e
encaminha solicitações de customização de produtos para grandes
compradores. Prepara documentação subsidiária para treinamento de
vendedores dos pontos de venda, em meio impresso e digital. Pode treinar
novos funcionários. Elabora relatórios e estatísticas. Comunica a sua
chefia os resultados de vendas. Consulta publicações técnicas, participa
de palestras e de feiras relacionadas à área. Participa de lançamento de
novos produtos.

\begin{Form}
\textbf{Esta correspondência está adequada?} \\ 
\ChoiceMenu[radio,radiosymbol=\ding{52},name=cbovalid_ 354130 ]{}{CBO deve ficar=sim, \hspace{6em} CBO não fica aqui=nao}
\end{Form}

\newpage
\section{ 05001 -- Técnico em Computação Gráfica }

\textbf{Perfil do Curso (CNCT)}

O Técnico em Computação Gráfica Digital será habilitado para:

Elaborar e implementar projetos de programação visual e lay-out para
mídia digital e/ou impressa. Realizar a modelagem e edição de imagens,
áudios e vídeos. Estruturar aplicações web e multimídia. Aplicar
técnicas de desenho e pintura digital. Realizar a programação de objetos
gráficos 2D e 3D. Realizar tratamento de imagens estáticas e em
movimento que compõem estruturas de navegação em mídias digitais.
Executar desenho técnico para elaboração de projetos, plantas e maquetes
digitais.

Para atuação como Técnico em Computação Gráfica, são fundamentais:

Conhecimentos e saberes relacionados aos processos de planejamento e
execução de projetos e roteiros de modo a garantir a entrega de produtos
digitais de acordo com suas finalidades. Conhecimentos e saberes
relacionados às normas técnicas, à liderança de equipes, à solução de
problemas técnicos e à assertividade na comunicação do material
produzido.

\textbf{CBOs não correspondidos e seus perfis:}

\textbf{CBO: 317120  ( Desenvolvedor de multimídia )}

Desenvolve sistemas e aplicações de multimídia, elaborando a interface
gráfica e aplicando critérios ergonômicos de navegação. Monta a
estrutura do banco de dados e codifica os programas e aplicativos,
aplicando sistemas de rotinas de segurança. Compila os programas. Testa
aplicativos e programas, elaborando casos de testes. Gera aplicativos
para instalação e gerenciamento dos sistemas. Documenta sistemas e
aplicações e avalia o desempenho dos produtos desenvolvidos. Implanta
sistemas e aplicações de multimídia, instalando os programas, adaptando
o conteúdo para mídias interativas e configurando os equipamentos em que
a aplicação será executada. Homologa sistemas e aplicações
desenvolvidos, avaliando e validando os resultados da implantação.
Elabora material para capacitação de usuários. Avalia os objetivos e as
metas de projetos de sistemas e aplicações, tendo em vista os resultados
obtidos na fase de implantação. Publica o código final no servidor de
sistemas e aplicações. Realiza a manutenção de sistemas e aplicações de
multimídia, atualizando informações gráficas, textuais e audiovisuais.
Implementa alterações em sistemas e aplicações e em suas estruturas de
armazenamento de dados, adequando-os aos sistemas e ambientes
operacionais. Monitora a performance de sistemas e aplicações,
utilizando ferramentas de software, para obter subsídios à atividade de
manutenção. Atualiza a documentação de sistemas e aplicações, em função
das intervenções de manutenção. Fornece suporte técnico para o cliente
interno. Colabora na seleção dos recursos de desenvolvimento,
especificando máquinas, ferramentas, acessórios e suprimentos. Participa
da seleção das linguagens de programação e metodologias. Seleciona
ferramentas de desenvolvimento. Participa da definição da nomenclatura
padrão. Colabora no planejamento das etapas e ações de trabalho,
participando das definições de atividades, tarefas e cronograma.
Participa de reuniões com a equipe de trabalho ou com o cliente.
Participa das definições de padronizações. Projeta sistemas e aplicações
de multimídia, identificando a demanda do cliente, coletando dados e
desenvolvendo o leiaute de telas e relatórios. Elabora o pré-projeto e
os projetos conceitual, lógico, estrutural, físico e gráfico dos
sistemas e aplicações, definindo os requisitos de projeto. Participa das
definições dos critérios ergonômicos de navegação e da interface de
comunicação e interatividade. Dimensiona a vida útil dos sistemas e
aplicações.

\begin{Form}
\textbf{Esta correspondência está adequada?} \\ 
\ChoiceMenu[radio,radiosymbol=\ding{52},name=cbovalid_ 317120 ]{}{CBO deve ficar=sim, \hspace{6em} CBO não fica aqui=nao}
\end{Form}

\newpage
\section{ 05003 -- Técnico em Informática }

\textbf{Perfil do Curso (CNCT)}

O Técnico em Informática será habilitado para:

Desenvolver sistemas computacionais utilizando ambiente de
desenvolvimento. Realizar modelagem, desenvolvimento, testes,
implementação e manutenção de sistemas computacionais. Modelar,
construir e realizar manutenção de banco de dados. Executar montagem,
instalação e configuração de equipamentos de informática. Instalar e
configurar sistemas operacionais e aplicativos em equipamentos
computacionais. Realizar manutenção preventiva e corretiva de
equipamentos de informática. Instalar e configurar dispositivos de
acesso à rede e realizar testes de conectividade. Realizar atendimento
help-desk. Operar, instalar, configurar e realizar manutenção em redes
de computadores. Aplicar técnicas de instalação e configuração da rede
física e lógica. Instalar, configurar e administrar sistemas
operacionais em redes de computadores. Executar as rotinas de
monitoramento do ambiente operacional. Identificar e registrar os
desvios e adotar os procedimentos de correção. Executar procedimentos de
segurança, pré-definidos, para ambiente de rede.

Para atuação como Técnico em Informática, são fundamentais:

Conhecimentos e saberes relacionados aos processos de planejamento e
execução de projetos computacionais de forma a garantir a entrega de
produtos digitais, análise de softwares, testagem de protótipos, de
acordo com suas finalidades. Conhecimentos e saberes relacionados às
normas técnicas, à liderança de equipes, à solução de problemas técnicos
e à assertividade na comunicação de laudos e análises. Habilidades
relacionadas à construção de soluções em BI e integrações sistêmicas.

\textbf{CBOs não correspondidos e seus perfis:}

\textbf{CBO: 317115  ( Programador de máquinas - ferramenta com comando numérico )}

Projeta sistemas e aplicações para máquinas-ferramenta CNC,
identificando a demanda do cliente, coletando dados e desenvolvendo o
leiaute de telas e relatórios. Elabora o pré-projeto e os projetos
conceitual, lógico, estrutural, físico e gráfico dos sistemas e
aplicações. Elabora croquis e desenhos para geração de programas de CNC.
Projeta dispositivos, ferramentas e posicionamento de peças em máquinas.
Participa das definições dos critérios ergonômicos de navegação e da
interface de comunicação e interatividade. Dimensiona a vida útil dos
sistemas e aplicações e modela sua estrutura de banco de dados.
Desenvolve sistemas e aplicações para máquinas-ferramenta CNC,
elaborando a interface gráfica, montando a estrutura do banco de dados,
codificando os programas e aplicativos e aplicando sistemas de rotinas
de segurança. Compila os programas. Testa aplicativos e programas,
elaborando casos de testes. Gera aplicativos para instalação e
gerenciamento dos sistemas. Documenta os sistemas e aplicações e avalia
o desempenho dos produtos desenvolvidos. Implanta sistemas e aplicações
para máquinas-ferramenta CNC, instalando os programas, homologando os
sistemas e aplicações desenvolvidos, avaliando e validando os resultados
da implantação. Elabora material para capacitação de usuários. Realiza a
manutenção de sistemas e aplicações para máquinas-ferramenta CNC,
atualizando informações gráficas, textuais e audiovisuais e convertendo
seus códigos fontes para outras linguagens ou plataformas. Implementa
alterações nos sistemas e aplicações. Monitora a performance dos
sistemas e aplicações, utilizando ferramentas de software, para obter
subsídios à atividade de manutenção. Atualiza a documentação dos
sistemas e aplicações, em função das intervenções de manutenção. Fornece
suporte técnico para o cliente interno. Colabora na seleção dos recursos
de desenvolvimento, especificando máquinas, ferramentas, acessórios e
suprimentos. Participa da seleção das linguagens de programação e
metodologias. Seleciona ferramentas de desenvolvimento. Colabora no
planejamento das etapas e ações de trabalho, participando das definições
de atividades, tarefas e cronograma. Participa de reuniões com a equipe
de trabalho ou com o cliente. Participa das definições de padronizações.

\begin{Form}
\textbf{Esta correspondência está adequada?} \\ 
\ChoiceMenu[radio,radiosymbol=\ding{52},name=cbovalid_ 317115 ]{}{CBO deve ficar=sim, \hspace{6em} CBO não fica aqui=nao}
\end{Form}

\textbf{CBO: 317120  ( Desenvolvedor de multimídia )}

Desenvolve sistemas e aplicações de multimídia, elaborando a interface
gráfica e aplicando critérios ergonômicos de navegação. Monta a
estrutura do banco de dados e codifica os programas e aplicativos,
aplicando sistemas de rotinas de segurança. Compila os programas. Testa
aplicativos e programas, elaborando casos de testes. Gera aplicativos
para instalação e gerenciamento dos sistemas. Documenta sistemas e
aplicações e avalia o desempenho dos produtos desenvolvidos. Implanta
sistemas e aplicações de multimídia, instalando os programas, adaptando
o conteúdo para mídias interativas e configurando os equipamentos em que
a aplicação será executada. Homologa sistemas e aplicações
desenvolvidos, avaliando e validando os resultados da implantação.
Elabora material para capacitação de usuários. Avalia os objetivos e as
metas de projetos de sistemas e aplicações, tendo em vista os resultados
obtidos na fase de implantação. Publica o código final no servidor de
sistemas e aplicações. Realiza a manutenção de sistemas e aplicações de
multimídia, atualizando informações gráficas, textuais e audiovisuais.
Implementa alterações em sistemas e aplicações e em suas estruturas de
armazenamento de dados, adequando-os aos sistemas e ambientes
operacionais. Monitora a performance de sistemas e aplicações,
utilizando ferramentas de software, para obter subsídios à atividade de
manutenção. Atualiza a documentação de sistemas e aplicações, em função
das intervenções de manutenção. Fornece suporte técnico para o cliente
interno. Colabora na seleção dos recursos de desenvolvimento,
especificando máquinas, ferramentas, acessórios e suprimentos. Participa
da seleção das linguagens de programação e metodologias. Seleciona
ferramentas de desenvolvimento. Participa da definição da nomenclatura
padrão. Colabora no planejamento das etapas e ações de trabalho,
participando das definições de atividades, tarefas e cronograma.
Participa de reuniões com a equipe de trabalho ou com o cliente.
Participa das definições de padronizações. Projeta sistemas e aplicações
de multimídia, identificando a demanda do cliente, coletando dados e
desenvolvendo o leiaute de telas e relatórios. Elabora o pré-projeto e
os projetos conceitual, lógico, estrutural, físico e gráfico dos
sistemas e aplicações, definindo os requisitos de projeto. Participa das
definições dos critérios ergonômicos de navegação e da interface de
comunicação e interatividade. Dimensiona a vida útil dos sistemas e
aplicações.

\begin{Form}
\textbf{Esta correspondência está adequada?} \\ 
\ChoiceMenu[radio,radiosymbol=\ding{52},name=cbovalid_ 317120 ]{}{CBO deve ficar=sim, \hspace{6em} CBO não fica aqui=nao}
\end{Form}

\newpage
\section{ 05009 -- Técnico em Telecomunicações }

\textbf{Perfil do Curso (CNCT)}

O Técnico em Telecomunicações será habilitado para:

Participar na elaboração de projetos de telecomunicações. Instalar,
testar e realizar manutenções preventivas e corretivas em sistemas de
telecomunicações. Configurar equipamentos nas áreas de telefonia,
transmissão e redes de comunicação. Supervisionar tecnicamente processos
e serviços de telecomunicações. Elaborar documentação técnica. Prestar
assistência técnica aos clientes. Realizar programação de softwares
específicos para equipamentos de telecomunicações. Participar na
elaboração da documentação técnica

Para atuação como Técnico em Telecomunicações, são fundamentais:

Conhecimentos e saberes relacionados aos processos técnicos de
telecomunicação cabeada ou de transmissão/tráfego de dados móveis, bem
como às boas práticas de comunicação e de liderança de equipes.

\textbf{CBOs não correspondidos e seus perfis:}

\textbf{CBO: 391205  ( Inspetor de qualidade )}

Planeja e desenvolve a inspeção de insumos e produtos - peças, conjuntos
e material em processo de produção -, definindo amostras e aplicando
procedimentos e métodos de inspeção de acordo com sistemas, normas e
instruções específicas. Define e opera equipamentos laboratoriais,
instrumentos e ferramentas para realização de ensaios. Pode fazer coleta
e envio de amostras para análises. Aplica ferramentas da qualidade e
metodologia estatística para avaliação da conformidade, de acordo com o
objeto da inspeção. Verifica conformidade de processos, controlando as
condições do ambiente em que se desenvolve o processo e observando
condições de uso de equipamentos, instrumentos e ferramentas. Analisa
relatórios de controle de processos para constatar o cumprimento de
normas e procedimentos e informar setores responsáveis sobre processos,
produtos e serviços não-conformes. Na inspeção de recebimento, na
inspeção realizada durante o processo de produção e na inspeção de
produtos acabados, examina documentos, verifica local de recebimento e
inspeção, avalia condições de descarregamento e verifica higiene e
limpeza do local de armazenamento de insumos e produtos. Pode executar
ações corretivas em função de reclamações do mercado. Libera produtos e
serviços inspecionados, verificando especificações e padrões e
conferindo condições de transporte e higiene.

\begin{Form}
\textbf{Esta correspondência está adequada?} \\ 
\ChoiceMenu[radio,radiosymbol=\ding{52},name=cbovalid_ 391205 ]{}{CBO deve ficar=sim, \hspace{6em} CBO não fica aqui=nao}
\end{Form}

\newpage
\section{ 06005 -- Técnico em Edificações }

\textbf{Perfil do Curso (CNCT)}

O Técnico em Edificações será habilitado para:

Desenvolver projetos de arquitetura, estrutura, instalações elétricas e
hidrossanitárias de até 80 m2 usando meios físicos ou digitais. Elaborar
orçamentos de obras e serviços. Planejar a execução dos serviços de
construção e manutenção predial. Executar obras e serviços de construção
e manutenção predial. Executar ensaios de materiais de construção, solos
e controle tecnológico. Conduzir planos de qualidade da construção.
Coordenar a execução de serviços de manutenção de equipamentos e
instalações em edificações.

Para atuação como Técnico em Edificações, são fundamentais:

Conhecimentos e saberes relacionados aos processos de planejamento e
construção de edificações de modo a assegurar a saúde e a segurança dos
trabalhadores e dos futuros ocupantes do imóvel. Conhecimentos e saberes
relacionados à sustentabilidade do processo produtivo, às técnicas e
processos de produção na construção civil, às normas técnicas.
Habilidades e competências relacionadas à liderança de equipes, à
solução de problemas técnicos e trabalhistas e à gestão de conflitos.

\textbf{CBOs não correspondidos e seus perfis:}

\textbf{CBO: 318005  ( Desenhista técnico )}

Analisa solicitação de desenho técnico, interpretando documentos de
apoio, tais como plantas, projetos, catálogos, croquis e normas. Propõe
ao solicitante alternativas para execução do desenho, chegando a um
acordo sobre os detalhes técnicos finais. Estima tempo para realização
do desenho e define prioridades, conforme cronograma. Programa ações
para realização do desenho técnico, observando características técnicas,
fazendo esboços, estabelecendo formatos e definindo escalas e sistemas
de representação. Elabora desenhos técnicos, representando formas
geométricas em duas e três dimensões. Utiliza computador, programa --
tal como CAD-Desenho Assistido por Computador (Computer Aided Design) -,
impressoras, plotters e outros periféricos, além de fazer uso de
instrumentos para execução do traçado diretamente em substratos de
papel. Constrói o desenho, traçando linhas auxiliares, cotando e
definindo a lista de materiais e componentes. Envia o desenho para
revisão, recebendo a aprovação do solicitante. Realiza cópias de
segurança e disponibiliza o desenho técnico final e/ou suas revisões
para as áreas afins, que providenciarão a entrega ao solicitante.
Conserva o local de trabalho limpo e organizado. Controla graus de
luminosidade do ambiente e verifica as condições ergonômicas para uso de
monitor de vídeo, mouse, mesa digitalizadora, prancheta e teclado.
Mantém equipamentos e instrumentos de trabalho em plenas condições de
uso e funcionamento. Requisita serviço de manutenção de computador,
impressoras, plotters e outros periféricos, quando necessário. Organiza
e arquiva documentos físicos e eletrônicos, tais como comprovantes de
aprovação e cópias de segurança dos desenhos técnicos. Destina
adequadamente os resíduos gerados, fazendo reciclagem e realizando
descarte de acordo com as normas ambientais.

\begin{Form}
\textbf{Esta correspondência está adequada?} \\ 
\ChoiceMenu[radio,radiosymbol=\ding{52},name=cbovalid_ 318005 ]{}{CBO deve ficar=sim, \hspace{6em} CBO não fica aqui=nao}
\end{Form}

\textbf{CBO: 318010  ( Desenhista copista )}

Analisa solicitação de cópia de desenho técnico, observando as
características técnicas do desenho original e interpretando documentos
de apoio, tais como originais de plantas, projetos, catálogos, croquis e
normas. Propõe ao solicitante alternativas para execução da cópia,
chegando a acordo em relação aos detalhes técnicos finais. Estima tempo
para realização do trabalho. Elabora cópias de desenhos técnicos,
reproduzindo representações de formas geométricas em duas e três
dimensões. Utiliza computador, programa -- tal como CAD-Desenho
Assistido por Computador (Computer Aided Design) -, mesas
digitalizadoras, impressoras, plotters e outros periféricos, além de
fazer uso de instrumentos para execução do traçado diretamente em
substratos de papel. Constrói a cópia do desenho, traçando linhas
auxiliares e cotando. Envia a cópia para revisão, recebendo a aprovação
do solicitante. Realiza cópias de segurança e disponibiliza o desenho
final copiado e/ou suas revisões para as áreas afins, que providenciarão
a entrega ao solicitante. Conserva o local de trabalho limpo e
organizado. Controla graus de luminosidade do ambiente e verifica as
condições ergonômicas para uso de monitor de vídeo, mouse, mesa
digitalizadora, prancheta e teclado. Mantém equipamentos e instrumentos
de trabalho em plenas condições de uso e funcionamento. Requisita
serviço de manutenção de computador, impressoras, plotters e outros
periféricos, quando necessário. Organiza e arquiva documentos físicos e
eletrônicos, tais como comprovantes de aprovação e uma via de cada
desenho copiado. Destina adequadamente os resíduos gerados, fazendo
reciclagem e realizando descarte de acordo com as normas ambientais.

\begin{Form}
\textbf{Esta correspondência está adequada?} \\ 
\ChoiceMenu[radio,radiosymbol=\ding{52},name=cbovalid_ 318010 ]{}{CBO deve ficar=sim, \hspace{6em} CBO não fica aqui=nao}
\end{Form}

\textbf{CBO: 318015  ( Desenhista detalhista )}

Analisa solicitação de detalhamento de desenho técnico, interpretando
original e documentos de apoio, tais como plantas, projetos, catálogos,
croquis e normas. Propõe ao solicitante alternativas para execução do
detalhamento, chegando a um acordo sobre pormenores técnicos finais.
Estima tempo para realização do detalhamento e estabelece prioridades,
conforme cronograma. Programa ações para realização do desenho
detalhado, observando características técnicas, fazendo esboços,
estabelecendo formatos e definindo escalas e sistemas de representação.
Elabora desenhos técnicos detalhados, representando formas geométricas
em duas e três dimensões. Utiliza computador, programa -- tal como
CAD-Desenho Assistido por Computador (Computer Aided Design) -,
impressoras, plotters e outros periféricos, além de fazer uso de
instrumentos para execução do traçado diretamente em substratos de
papel. Constrói o detalhamento do desenho, traçando linhas auxiliares,
cotando e realçando aspectos e características de formas geométricas
complexas e outros elementos a detalhar. Elabora a lista de materiais e
componentes do desenho detalhado. Envia o desenho técnico detalhado para
revisão, recebendo a aprovação do solicitante. Realiza cópias de
segurança e disponibiliza o desenho detalhado final e/ou suas revisões
para as áreas afins, que providenciarão a entrega ao solicitante.
Conserva o local de trabalho limpo e organizado. Controla graus de
luminosidade do ambiente e verifica as condições ergonômicas para uso de
monitor de vídeo, mouse, mesa digitalizadora, prancheta e teclado.
Mantém equipamentos e instrumentos de trabalho em plenas condições de
uso e funcionamento. Requisita serviço de manutenção de computador,
impressoras, plotters e outros periféricos, quando necessário. Organiza
e arquiva documentos físicos e eletrônicos, tais como comprovantes de
aprovação e cópias de segurança dos desenhos técnicos detalhados.
Destina adequadamente os resíduos gerados, fazendo reciclagem e
realizando descarte de acordo com as normas ambientais.

\begin{Form}
\textbf{Esta correspondência está adequada?} \\ 
\ChoiceMenu[radio,radiosymbol=\ding{52},name=cbovalid_ 318015 ]{}{CBO deve ficar=sim, \hspace{6em} CBO não fica aqui=nao}
\end{Form}

\textbf{CBO: 318120  ( Desenhista técnico (instalações hidrossanitárias) )}

Planeja o trabalho relacionado ao desenho técnico de instalações
hidrossanitárias, propondo alternativas para elaboração do desenho e
determinando a metodologia para sua execução. Seleciona instrumentos e
outros meios necessários à elaboração do desenho. Pode prever uso de
sistemas BIM - Modelagem da Informação da Construção (Building
Information Modeling). Calcula e define o custo do desenho. Estipula
prazo de execução. Prepara o local de trabalho, de acordo com
recomendações normativas de ergonomia, de segurança do trabalho e de
saúde ocupacional. Coleta dados para elaboração do desenho técnico de
instalações hidrossanitárias, interpretando projetos existentes,
consultando informações em arquivos e fazendo levantamento de campo.
Busca informações complementares e organiza os dados coletados. Processa
dados para desenho técnico de instalações hidrossanitárias, analisando
croquis obtidos por meio dos levantamentos de campo; interpretando
memória de cálculo e outras informações do conjunto de dados coletados.
Transfere dados da estação local de campo -- para o escritório ou
ambiente onde será executado o desenho --, utilizando ferramentas de
tecnologia da informação requeridas. Calcula dados da caderneta de
campo. Realiza pesquisas complementares -- inclusive via internet --,
para obter novos dados relacionados ao desenho ou ampliar a compreensão
sobre o conjunto de dados à disposição. Elabora desenhos técnicos de
instalações hidrossanitárias -- tais como desenhos de redes de esgoto,
águas pluviais, abastecimento de água e de sistemas de captação,
purificação e distribuição de água --, aplicando normas técnicas e
realizando visitas técnicas para rever dados. Utiliza softwares
específicos para desenho, especialmente CAD-Desenho Assistido por
Computador (Computer Aided Design). Utiliza estações de trabalho
computadorizadas, impressoras, plotters e outros periféricos. Pode fazer
uso de instrumentos para execução do traçado diretamente em substratos
de papel, quando requerido. Define formatos e escalas e detalha
desenhos. Diagrama pranchas ou folhas de desenho técnico. Legenda
plantas e elabora a lista de material relacionada ao conteúdo do
desenho. Revisa desenhos técnicos de instalações hidrossanitárias,
conferindo cotas, dimensionamentos, informações descritivas e outros
conteúdos. Atualiza desenhos, em função de mudanças de legislação.
Realiza a complementação ou a modificação de desenhos, quando
necessário. Em situações em que os desenhos passam por supervisão,
realiza as compatibilizações necessárias dos desenhos, de acordo com os
apontamentos da supervisão. Fecha a ordem de serviço do desenho de
instalações hidrossanitárias, relatando mudanças de procedimento
adotadas ao longo dos trabalhos; liberando o desenho para arquivo
eletrônico ou mapoteca; e enviando originais e/ou cópias do desenho para
clientes. Retorna documentos utilizados para arquivo. Organiza arquivos
técnicos -- físicos e eletrônicos -- relacionados aos serviços com
desenhos técnicos de instalações hidrossanitárias, determinando o tipo
de arquivo a ser utilizado e indexando documentos pertinentes à área.
Reúne documentos e organiza catálogos de fornecedores. Compacta e
armazena arquivos digitais. Conserva o local de trabalho limpo e
organizado. Controla luminosidade do ambiente e verifica as condições
ergonômicas para uso de monitor de vídeo, mouse, mesa digitalizadora,
prancheta e teclado. Mantém equipamentos e instrumentos de trabalho em
plenas condições de uso e funcionamento. Requisita serviço de manutenção
de computador, impressoras, plotters e outros periféricos, quando
necessário. Destina adequadamente os resíduos gerados, fazendo
reciclagem e realizando descarte de acordo com as normas ambientais.
Zela pela segurança, atuando na prevenção de acidentes e utilizando
equipamentos de proteção individual, sempre que necessário,
especialmente durante os levantamentos de dados em campo.

\begin{Form}
\textbf{Esta correspondência está adequada?} \\ 
\ChoiceMenu[radio,radiosymbol=\ding{52},name=cbovalid_ 318120 ]{}{CBO deve ficar=sim, \hspace{6em} CBO não fica aqui=nao}
\end{Form}

\newpage
\section{ 06011 -- Técnico em Portos }

\textbf{Perfil do Curso (CNCT)}

O Técnico em Portos será habilitado para:

Desenvolver atividades de gerenciamento, monitoramento, supervisão,
programação e controle em operações portuárias diversas. Controlar,
programar e coordenar operações de transportes em geral, inclusive o
transporte de cargas perigosas. Prestar suporte técnico em atividades de
armazenagem de cargas, inclusive armazenagem de cargas perigosas.
Supervisionar operações de embarque, transbordo e desembarque de cargas
entre os diversos modos de transporte. Prestar suporte técnico para o
agenciamento de embarcações. Encaminhar procedimentos de importação e
exportação. Verificar as condições de segurança dos meios de
transportes, equipamentos utilizados e das cargas. Programar e
supervisionar a manutenção de equipamentos eletromecânicos de operação
portuária. Verificar e inspecionar a eficiência operacional de
equipamentos e veículos. Interpretar, elaborar e preparar a documentação
necessária ao desembaraço aduaneiro de cargas. Atender clientes internos
e externos. Elaborar a cotação de preços de serviços de transporte,
inclusive transporte multimodal. Identificar e programar rotas de
transporte de cargas. Utilizar tecnologias aplicadas ao processo de
gestão da informação sobre condições do transporte e da carga.

Para atuação como Técnico em Portos, são fundamentais:

Conhecimentos e saberes relacionados aos processos de desembaraço
aduaneiro de cargas, transporte terrestre de contêineres, operações
logísticas, transporte e armazenagem de mercadorias perigosas,
sistemáticas de importação e exportação, operações de
embarque/desembarque de navios e logística de armazéns. Comprometimento
com as questões ambientais, sociais e de desenvolvimento tecnológico e
para a solução de problemas e busca por inovações tecnológicas.
Conhecimentos relacionados à liderança de equipes, à solução de
problemas técnicos e trabalhistas e à gestão de conflitos.

\textbf{CBOs não correspondidos e seus perfis:}

\textbf{CBO: 354305  ( Analista de exportação e importação )}

Analisa mercado internacional de produtos e serviços, para identificar
oportunidades de negócio e fornecer subsídios para definição de
políticas de importação ou exportação. Examina aspectos cambiais e
creditícios. Analisa aspectos legais e normativos. Verifica as
tendências de consumo nos países, para adequar produtos e embalagens aos
diferentes mercados. Analisa logística de movimentação. Elabora o plano
de exportação, estabelecendo tipos e quantidades de mercadorias e
avaliando concorrentes. Participa da definição de preços de produtos e
serviços. Avalia a capacidade produtiva da mercadoria a exportar, para
projetar curvas de vendas e definir logística. Operacionaliza a
exportação, negociando condições de venda de produtos e serviços.
Elabora proposta comercial, incluindo especificações técnicas do produto
e tipos de embalagens e acondicionamento. Comunica ao produtor a
aceitação da proposta, acertando a compatibilização do cronograma de
produção com o contrato de venda. Contrata despachante, elabora
documentação para exportação e providencia pagamento de encargos.
Acompanha o andamento da produção de bens e serviços, para liberação da
mercadoria comprometida em tempo hábil de embarque. Contrata serviços de
terceiros para realização de transporte da mercadoria. Providencia meios
do recebimento do valor acordado, por intermédio de transações em moeda
estrangeira. Conclui o processo junto aos órgãos competentes,
acompanhando a finalização no Sistema Integrado de Comércio Exterior --
SISCOMEX. Guarda a documentação pelo prazo legal. Processa a importação
de mercadorias, estudando a viabilidade do processo e emitindo planilha
de custos. Verifica a conformidade da documentação e da nomenclatura das
mercadorias, em relação à legislação vigente. Faz o pedido das
mercadorias. No recebimento, instrui o desembaraço aduaneiro e verifica
a documentação. Pode negociar com fornecedores a troca de mercadorias
com defeito ou em desacordo com o pedido. Providencia meios de pagamento
por intermédio de transações em moeda estrangeira. Aciona a seguradora,
quando necessário. Realiza atividades de promoção de comércio exterior,
organizando missões empresariais, prestando assistência técnica e
operacional aos clientes em feiras e eventos e providenciando propaganda
de produtos e serviços. Atende consultas de clientes, orienta
exportadores e importadores e presta informações sobre o andamento de
processos de importação ou exportação. Desenvolve fornecedores. Pode
ministrar palestras técnicas da área.

\begin{Form}
\textbf{Esta correspondência está adequada?} \\ 
\ChoiceMenu[radio,radiosymbol=\ding{52},name=cbovalid_ 354305 ]{}{CBO deve ficar=sim, \hspace{6em} CBO não fica aqui=nao}
\end{Form}

\newpage
\section{ 06017 -- Técnico em Transporte Metroferroviário }

\textbf{Perfil do Curso (CNCT)}

O Técnico em Transporte Metroferroviário será habilitado para:

Operar, coordenar e fiscalizar o transporte metroferroviário. Coordenar
a circulação de veículos metroferroviários. Trabalhar no planejamento,
na execução e no controle de atividades ligadas às operações de pátios e
terminais, veículos, sinalização e equipamentos do transporte
metroferroviário. Coordenar a circulação de veículos metroferroviários.
Administrar e controlar as atividades de pátios e terminais. Controlar e
programar os horários de circulação de trens. Operar equipamentos e
sistemas de sinalização, telecomunicações e bordo utilizados nos
sistemas metroferroviários. Manobrar equipamentos e veículos
metroferroviários. Preencher relatórios, planilhas, documentos de
despacho, diário operacional e boletins de ocorrência.

Para atuação como Técnico em Transporte Metroferroviário, são
fundamentais:

Conhecimentos e saberes relacionados ao controle e operação de veículos
em pátios e vias ferroviárias, de modo a assegurar a saúde e a segurança
dos trabalhadores e dos usuários do transporte metroferroviário. Além de
ter o comprometimento com as questões sociais e de desenvolvimento
tecnológico, proatividade, formação para a solução de problemas e a
busca de inovações tecnológicas. Habilidades e saberes relacionados à
liderança de equipes, à solução de problemas técnicos e trabalhistas e à
gestão de conflitos.

\textbf{CBOs não correspondidos e seus perfis:}

\textbf{CBO: 342120  ( Afretador )}

Atende clientes, identificando suas necessidades. Informa-se sobre
produto a ser transportado, elaborando proposta de preços e serviços.
Define parcerias e fornecedores de serviços de armazenagem e transporte.
Contrata serviços de arrendamento de navio em diferentes modalidades,
entre as quais a casco nu (ou demise charter party ), afretamento por
tempo (time charter-party) e contrato por viagem (ou voyage party).
Contrata serviços de arrendamento de equipamentos, negociando preços e
prazos de pagamento. Elabora ou auxilia a elaboração de documentos e
contratos de afretamento. Registra dados de carga e transporte no
sistema de comércio exterior (Siscarga) e em sistemas internos. Confere,
organiza e arquiva documentos. Disponibiliza informações da carga para o
cliente, verificando a sua satisfação. Pesquisa mercado, coletando dados
de clientes potenciais. Identifica rotas de transporte, preços e
serviços de seus parceiros e de seus concorrentes. Consulta sistemas de
informação especializados, órgãos governamentais e entidades.

\begin{Form}
\textbf{Esta correspondência está adequada?} \\ 
\ChoiceMenu[radio,radiosymbol=\ding{52},name=cbovalid_ 342120 ]{}{CBO deve ficar=sim, \hspace{6em} CBO não fica aqui=nao}
\end{Form}

\textbf{CBO: 342315  ( Supervisor de carga e descarga )}

Supervisiona operações logísticas de carga e descarga, envolvendo o
embarque, o percurso e o desembarque. Verifica o tipo e a quantidade de
carga, emitindo documento fiscal. Examina as condições do transporte,
determinando as providências para adaptá-lo a cada tipo de embalagem de
carga e para adequá-lo às operações em terminais e armazéns. Confere a
retirada da carga. Colabora no planejamento, na programação e no
controle de atividades de transporte relacionadas a operações logísticas
de carga e descarga, identificando fatores condicionantes como limites
de horários para o transporte urbano de carga - principalmente nas
grandes cidades -, condições estabelecidas para movimentação de
materiais e determinação de formas de distribuição. Consulta sistemas de
informações, para verificar ocorrências no transporte de cargas.
Controla despesas, conferindo notas fiscais e prestações de contas.
Atende clientes, prestando informações.

\begin{Form}
\textbf{Esta correspondência está adequada?} \\ 
\ChoiceMenu[radio,radiosymbol=\ding{52},name=cbovalid_ 342315 ]{}{CBO deve ficar=sim, \hspace{6em} CBO não fica aqui=nao}
\end{Form}

\newpage
\section{ 06018 -- Técnico em Transporte Rodoviário }

\textbf{Perfil do Curso (CNCT)}

O Técnico em Transporte Rodoviário será habilitado para:

Organizar e controlar as operações e executar a logística de tráfego
rodoviário. Planejar, operacionalizar e executar a logística do
transporte de passageiros. Administrar e controlar a frota de veículos
no transporte rodoviário de cargas e passageiros. Executar a operação,
comercialização e manutenção de equipamentos. Planejar a armazenagem e o
processo de expedição das empresas e centros de distribuições. Planejar
e executar a distribuição de pessoal e cargas. Coordenar ações de
intermodalidade de transportes. Identificar as características da malha
viária e os diversos tipos de veículos transportadores. Aplicar a
legislação de trânsito de veículos e de transporte de passageiros.
Preparar e gerenciar a documentação necessária para operações de
transportes.

Para atuação como Técnico em Transporte Rodoviário, são fundamentais:

Conhecimentos e saberes relacionados ao controle de operações e
logística de transportes rodoviários, de passageiros e de cargas, de
modo a assegurar a saúde e a segurança dos trabalhadores e usuários.
Compromisso e ética com as questões ambientais, sociais e de
desenvolvimento tecnológico. Habilidades para a solução de problemas.
Conhecimentos e habilidades relacionados à busca de inovações
tecnológicas, à liderança de equipes, à solução de problemas técnicos e
trabalhistas e à gestão de conflitos.

\textbf{CBOs não correspondidos e seus perfis:}

\textbf{CBO: 342210  ( Despachante aduaneiro )}

Executa atividades relacionadas ao despacho aduaneiro de mercadorias -
inclusive bagagem de viajante - na importação, na exportação e na
internação, realizadas por qualquer via de transporte. Elabora e/ou
supervisiona a elaboração de documentos relativos ao despacho aduaneiro.
Providencia credenciamento de cliente junto à Receita Federal. Analisa
cartas de crédito. Elabora fatura proforma e fatura comercial. Prepara
lista de embarque. Executa e/ou supervisiona a classificação de
mercadorias, analisando a documentação técnica das mercadorias;
verificando suas funções de uso e aplicação; e observando seu material
constitutivo. Analisa necessidade de anuência de órgãos pertinentes e,
quando necessário, consulta a Receita Federal sobre dúvidas de
classificação. Analisa enquadramento legal das alíquotas de impostos de
importação e exportação. Enquadra mercadorias na Tarifa Externa Comum
(TEC), na Nomenclatura da Associação Latino-americana de Integração
(Naladi) e na Tabela de Incidência do Imposto sobre Produtos
Industrializados (TIPI). Preenche Certificado de Origem; Conhecimento de
Transporte internacional; Documento de Arrecadação de Receitas Federais
(DARF); guia de Imposto sobre Operações relativas à Circulação de
Mercadorias e sobre Prestações de Serviços de Transporte Interestadual e
Intermunicipal e de Comunicação (ICMS); além de outros formulários
necessários às operações de comércio exterior. Realiza e/ou supervisiona
a atividade de entrada de dados no Sistema Integrado de Comércio
Exterior (Siscomex) e sistemas auxiliares. Coleta dados de operações e
registra licenças e declarações de importação e de despacho de
exportação. Efetua registros no módulo Registro de Operações Financeiras
do Sistema Registro Declaratório Eletrônico (RDE-ROF) e de atos
concessórios de regime aduaneiro especial (``Drawback''). Faz Declaração
de Trânsito Aduaneiro (DTA). Efetua acompanhamento da tramitação de
documentos relacionados ao despacho aduaneiro. Solicita deferimento de
licenças de importação e de declarações retificadoras. Subscreve
documentos, inclusive termos de responsabilidade. Faz recebimento de
intimações, de notificações, de autos de infração, de despachos, de
decisões e de outros atos e termos processuais relacionados com o
procedimento de despacho aduaneiro. Realiza ou supervisiona o
recebimento de mercadorias e bagagens. Paga impostos e taxas. Pode
apresentar desistência de vistoria aduaneira. Executa acompanhamento de
vistoria da mercadoria na conferência aduaneira, observando a
conferência física e a retirada de amostra de mercadoria para
assistência técnica e perícia. Avalia exigências da receita federal e
demais órgãos pertinentes, presta esclarecimentos e apresenta
documentação. Retira comprovantes de importação e exportação junto à
receita federal. Presta assessoria a importadores e exportadores,
identificando suas necessidades; informando sobre legislação pertinente;
propondo alternativas de logística de transportes; e avaliando
benefícios fiscais. Calcula custos de operação. Pleiteia benefícios
fiscais. Monitora processos administrativos e informa a cliente sobre
tramitação de processos. Realiza contratação de serviços de terceiros,
supervisionando a cotação de preços de armazenagem e preços de fretes
nacionais e internacionais; orçando seguros; e negociando preços e
prazos. Confere cobrança de terceiros por serviços prestados e
providencia o pagamento. Supervisiona equipe de trabalho, distribuindo
tarefas, avaliando desempenho e desenvolvendo treinamentos. Monitora a
limpeza e a organização do ambiente de trabalho. Controla a aplicação de
procedimentos da empresa para cumprimento de normas ambientais,
inspecionando o reaproveitamento de sobras de material e o descarte de
resíduos.

\begin{Form}
\textbf{Esta correspondência está adequada?} \\ 
\ChoiceMenu[radio,radiosymbol=\ding{52},name=cbovalid_ 342210 ]{}{CBO deve ficar=sim, \hspace{6em} CBO não fica aqui=nao}
\end{Form}

\newpage
\section{ 09001 -- Técnico em Artes Circenses }

\textbf{Perfil do Curso (CNCT)}

O Técnico em Artes Circenses será habilitado para:

Criar, desenvolver e executar apresentações circenses em espaços de
circo, teatro, estúdio de televisão, públicos e culturais. Utilizar
técnicas artísticas e corporais de acrobacia aérea e de solo,
equilibrismo, malabarismo, antipodismo, ilusionismo, comicidade, canto,
dança e pantomima. Organizar e supervisionar a estrutura, a montagem e o
funcionamento do circo e dos equipamentos. Administrar, produzir e
divulgar espetáculos.

Para atuação como Técnico em Artes Circenses, são fundamentais:

Conhecimentos interdisciplinares relacionados aos processos de criação,
envolvendo pesquisa, idealização, planejamento, execução técnica,
fruição e recepção estética. Competências comunicativas e empreendedoras
voltadas à proposição de projetos, ao coletivo, à gestão, à solução de
problemas e à resiliência, entre outras competências socioemocionais.

\textbf{CBOs não correspondidos e seus perfis:}

\textbf{CBO: 376305  ( Apresentador de eventos )}

Prepara-se para realizar apresentação de evento, pesquisando informações
a transmitir na ocasião. Levanta informações impressas, radiofônicas,
televisivas e disponíveis na Internet. Contata fontes de informação -
tais como líderes comunitários, formadores de opinião e políticos --
para inteirar-se dos temas relacionados ao evento. Pesquisa o perfil
social, econômico e político do público-alvo. Atenta para
especificidades culturais e políticas regionais. Com base nas
informações coletadas, prepara comentários para fazer durante o evento e
capacita-se a entrevistar convidados. Verifica as condições técnicas
exigidas para a apresentação. Pode gravar locuções e materiais
audiovisuais a serem utilizados no evento como complementares ao
processo de comunicação. Prepara-se física e mentalmente para realizar a
apresentação de evento, selecionando vestimenta e adereços -- entre os
considerados adequados pela organização -; ajustando postura; realizando
exercícios de concentração; e aquecendo a voz. Pode interagir com os
membros da equipe de produção e organização, opinando na formatação do
evento, quando solicitado. Apresenta o evento, dando o tom da solenidade
e interpretando o texto. Pode receber instruções dos organizadores do
evento, durante sua apresentação. Conduz o evento, atento ao roteiro
predefinido. Alerta para colocação e retirada dos equipamentos.
Administra os tempos do evento. Presta informações ao público,
apresentando o conteúdo do evento. Anima o evento, mantendo a atenção e
criando expectativas no público. Controla o tempo das performances de
artistas e indica o momento de apresentação dos materiais audiovisuais
previamente gravados. Entrevista convidados, aprofundando as informações
por eles trazidas e coordenando a apresentação de dúvidas pelo público e
de respostas pelos entrevistados. Transmite credibilidade para o
público, noticiando fatos e imprevistos de forma segura. Faz adequação
de sua postura física e entonação de voz à informação a ser transmitida.
Verifica o conteúdo e a forma do texto, para atender às expectativas do
público. Observa sinais de entrada de inserções comerciais. Transmite
informações sobre o patrocinador, divulgando o produto promocional.
Informa o público sobre acontecimentos intervenientes ao evento, tais
como alertas sobre criança perdida, trânsito, entre outros. Percebe as
reações relacionadas ao nível de satisfação do público durante o evento
e conversa com alguns participantes após o final, para conhecer suas
opiniões. Pode elaborar relatório para as equipes de produção e
organização, com dados sobre ocorrências durante o evento e com
sugestões para os próximos. Organiza e arquiva - em meio físico e
eletrônico -- os dados coletados, para reutilização em eventos com temas
similares.

\begin{Form}
\textbf{Esta correspondência está adequada?} \\ 
\ChoiceMenu[radio,radiosymbol=\ding{52},name=cbovalid_ 376305 ]{}{CBO deve ficar=sim, \hspace{6em} CBO não fica aqui=nao}
\end{Form}

\textbf{CBO: 376310  ( Apresentador de festas populares )}

Prepara-se para realizar apresentação de festa popular, pesquisando
informações a transmitir na ocasião. Levanta informações impressas,
radiofônicas, televisivas e disponíveis na Internet sobre artistas
participantes. Contata fontes de informação - tais como líderes
comunitários e formadores de opinião -- para inteirar-se de temas
relacionados à festa. Pesquisa o perfil socioeconômico e a faixa etária
do público-alvo. Atenta para especificidades -- tais como região de
origem, ano de início, tradições culinárias, danças, músicas,
brincadeiras realizadas, entre outras - relacionadas à festa. Com base
nas informações coletadas, prepara comentários para fazer durante a
festa e capacita-se a entrevistar convidados. Verifica as condições
técnicas exigidas para a apresentação. Pode gravar locuções e materiais
audiovisuais a serem utilizados na festa como complementares ao processo
de comunicação. Prepara-se física e mentalmente para realizar a
apresentação da festa, selecionando vestimenta e adereços -- entre os
considerados adequados pela organização -; ajustando postura; realizando
exercícios de concentração; e aquecendo a voz. Pode interagir com os
membros da equipe de produção e organização, opinando na formatação da
festa, quando solicitado. Apresenta a festa, mantendo a espontaneidade e
controlando a gesticulação. Pode receber instruções dos organizadores da
festa, durante sua apresentação. Conduz a festa, atento ao roteiro
predefinido. Alerta para colocação e retirada de equipamentos e de
instrumentos musicais. Administra o tempo para realização de cada item
do roteiro. Interage com o público, estabelecendo relações de empatia e
estimulando o convívio cordial entre os participantes. Anima a festa,
mantendo a atenção e criando expectativas no público. Controla o tempo
da performance de cada artista e indica o momento de apresentação dos
materiais audiovisuais previamente gravados. Entrevista convidados,
aprofundando as informações por eles trazidas e sendo porta-voz de
dúvidas e indagações do público. Transmite informações e comunicados
relevantes ao desenvolvimento da festividade. Faz adequação de sua
postura física e entonação de voz à informação a ser transmitida.
Verifica os conteúdos de textos e a forma de transmiti-los, para não
dispersar a atenção do público. Informa ao público sobre ocorrências
durante a festa, tais como alertas sobre criança perdida, estacionamento
irregular de veículos, entre outras. Faz contato - por meio do olhar -
com o público, para perceber suas reações durante a festa e levantar o
seu nível de satisfação. Conversa com alguns participantes após o final
da festa, para conhecer suas opiniões. Pode elaborar relatório para as
equipes de produção e organização, com dados sobre ocorrências durante a
festa e com sugestões para as próximas. Organiza e arquiva - em meio
físico e eletrônico -- os dados coletados sobre a festa, que possam ser
reutilizados em novas edições, nos anos seguintes.

\begin{Form}
\textbf{Esta correspondência está adequada?} \\ 
\ChoiceMenu[radio,radiosymbol=\ding{52},name=cbovalid_ 376310 ]{}{CBO deve ficar=sim, \hspace{6em} CBO não fica aqui=nao}
\end{Form}

\textbf{CBO: 376315  ( Apresentador de programas de rádio )}

Prepara-se para realizar apresentação de programa de rádio, levantando o
conteúdo básico de seu programa com a equipe de produção e pesquisando
informações a transmitir. Pesquisa características do seu padrão de
ouvinte ou público-alvo, tais como perfil socioeconômico, faixa etária,
uso de celular e Internet, principais tipos de divertimento ou
entretenimento; principais problemas enfrentados no dia a dia, entre
outras. A partir do conteúdo central de seu programa, levanta -- em
jornais impressos e suas versões na web; em programas jornalísticos em
rádios ou televisão; em dados disponíveis na Internet; e em outros meios
de comunicação - informações sobre diferentes temas, para manter-se
atualizado e poder explicar o assunto com clareza para seu público.
Examina os temas de maior interesse para seu público-alvo e os mais
discutidos em redes sociais. Contata fontes - tais como líderes
comunitários, formadores de opinião e políticos --, para reunir novas
informações. Com base nas informações coletadas, prepara comentários e
argumentos para seus programas e capacita-se a entrevistar convidados.
Verifica as condições técnicas exigidas para a apresentação. Pode
realizar programas de rádio ao vivo ou gravá-los, considerando todas as
etapas de produção e as expectativas do público-alvo. Pode transmitir
evento ao vivo no programa de rádio, atuando remotamente. Pode interagir
com os membros da equipe de produção da rádio, opinando sobre pauta,
gênero e formato de programa, quando solicitado. Prepara-se para
realizar a apresentação do programa de rádio, aquecendo sua voz e
controlando sua respiração. Busca adequar sua linguagem ao público-alvo,
para tornar clara sua fala. Apresenta o programa, interpretando textos.
Pode receber instruções da equipe de produção, durante sua apresentação.
Conduz o programa, atento ao roteiro predefinido. Administra o tempo
para apresentar cada item do roteiro. Informa o público, apresentando o
conteúdo do programa. Anima o programa, mantendo a atenção e criando
expectativas no público. Conversa com artistas e apresenta as músicas
preferidas de seu público-alvo. Interage com repórter, para detalhar
ocorrências do momento, ficando atento ao tempo da reportagem. Recebe
especialistas em temas abordados, para comentar notícias. Recebe
convidados, para realizar entrevistas e aprofundar as informações por
eles trazidas. Interage com os ouvintes, levantando problemas
enfrentados e reivindicando respostas e soluções de autoridades.
Transmite credibilidade para o público, noticiando fatos e imprevistos
de forma segura. Faz adequação da entonação de sua voz à informação a
ser transmitida. Observa sinais de entrada de inserções comerciais.
Transmite informações sobre o patrocinador e divulga produto
promocional. Recebe relatórios de audiência. Pode acompanhar reunião de
``grupo de foco'', que compartilha suas opiniões, ideias e reações em
relação ao programa. Pode conversar com amostra de ouvintes, para
conhecer suas opiniões e levantar pontos considerados positivos e
negativos. Com essas informações, pode ser chamado, pela equipe de
produção, para participar da discussão de alterações no formato do
programa e ajustes nos conteúdos abordados, para aumentar a audiência.
Organiza e arquiva - em meio físico e eletrônico -- as informações
coletadas em suas pesquisas, para consultas.

\begin{Form}
\textbf{Esta correspondência está adequada?} \\ 
\ChoiceMenu[radio,radiosymbol=\ding{52},name=cbovalid_ 376315 ]{}{CBO deve ficar=sim, \hspace{6em} CBO não fica aqui=nao}
\end{Form}

\textbf{CBO: 376320  ( Apresentador de programas de televisão )}

Prepara-se para realizar apresentação de programa de televisão,
consultando o conteúdo básico de seu programa com equipe de produção e
levantando informações. Pesquisa características do seu público - tais
como perfil socioeconômico, gênero e faixa etária -, para antever
expectativas e interesses. Mantém-se atualizado sobre diferentes temas
-- cultura, esporte, economia, política, entre outros -- do país e do
mundo, lendo livros, estudando artigos e matérias publicadas e
verificando informações divulgadas em jornais impressos e suas versões
na web; em programas jornalísticos de rádio e televisão; na Internet; e
em outros meios de comunicação. Examina os temas de maior interesse para
seu público e os mais discutidos em redes sociais. Contata fontes - tais
como líderes comunitários, formadores de opinião e políticos --, para
levantar informações locais ou regionais. Com base nas informações
coletadas, prepara comentários e argumentos para seus programas e
capacita-se a entrevistar convidados. Verifica as condições técnicas
exigidas para a apresentação. Pode interagir com diretor do programa e
com equipe de produção, opinando sobre pauta e formatação do programa,
quando solicitado. Mantém bom relacionamento com a equipe do programa,
trocando ideias. Prepara-se para realizar a apresentação do programa de
televisão, selecionando vestimenta e adereços -- entre os considerados
adequados pela direção -; ajustando postura; realizando exercícios de
concentração; aquecendo a voz; e controlando respiração. Busca adequar
sua linguagem ao seu público, para tornar clara sua fala. Usa
concentração e disciplina para eliminar vícios de linguagem. Apresenta
programa em uma área específica -- como culinária, esporte, música,
religião, entre outras -- ou com pauta diversificada; com ou sem
plateia; ao vivo ou gravado e editado. Atua em estúdio ou remotamente,
principalmente em transmissão de eventos ao vivo. Conduz o programa,
administrando o tempo para apresentar cada item previsto em roteiro.
Pode receber instruções do diretor ou da equipe de produção, durante sua
apresentação. Anima o programa, mantendo a atenção e criando
expectativas no público; demonstrando espontaneidade e descontração; e
desenvolvendo formas de comunicação para conquistar a empatia das
pessoas. Lida com imprevistos, durante a apresentação, de forma
criativa. Busca ter estilo pessoal de condução do programa,
diferenciando-se dos demais apresentadores. Informa o público,
apresentando o conteúdo do programa. Interage com os telespectadores,
levantando problemas enfrentados e reivindicando respostas e soluções de
autoridades. Conversa com artistas e apresenta as músicas preferidas de
seu público. Dialoga com especialistas em temas abordados, para comentar
notícias. Recebe convidados, para realizar entrevistas e aprofundar as
informações por eles trazidas. No caso de telejornais, segue o
``espelho'' do jornal (tabela com o tempo de cada matéria e a ordem de
apresentação). Interpreta textos que aparecem no ``teleprompter'', que
fica na câmera para a qual se dirige. Apresenta matérias gravadas e
editadas. Interage com repórter em tempo real, para detalhar ocorrências
do momento. Fica atento às orientações do diretor de TV. Transmite
credibilidade, noticiando fatos e imprevistos de forma segura. Faz
adequação da entonação de sua voz à informação a ser transmitida.
Observa sinais de entrada de inserções comerciais. Pode transmitir
informações sobre o patrocinador e divulgar produtos. Recebe relatórios
de audiência. Analisa pontos fortes e fracos de seu programa, levantados
por diferentes estratégias, como reunião de ``grupo de foco'' e coleta
de dados junto ao público. Com essas informações, pode ser chamado, pelo
diretor de TV e pela equipe de produção, para participar da discussão de
alterações no formato do programa e ajustes nos conteúdos abordados,
para aumentar a audiência.

\begin{Form}
\textbf{Esta correspondência está adequada?} \\ 
\ChoiceMenu[radio,radiosymbol=\ding{52},name=cbovalid_ 376320 ]{}{CBO deve ficar=sim, \hspace{6em} CBO não fica aqui=nao}
\end{Form}

\newpage
\section{ 09011 -- Técnico em Design de Interiores }

\textbf{Perfil do Curso (CNCT)}

O Técnico em Design de Interiores será habilitado para:

Criar, desenvolver e viabilizar a execução de projetos de interiores
residenciais, comerciais, de vitrines e exposições. Orientar e
desenvolver projetos com base em ergonomia e desenho universal,
conforto, saúde e bem-estar. Desenvolver esboços, perspectivas e
desenhos técnicos. Planejar e organizar o espaço, com base nos estudos
ergonômicos, estéticos e funcionais. Identificar elementos básicos para
a concepção projetual dos espaços internos habitados. Representar os
elementos de projeto no espaço bi e tridimensional. Aplicar métodos de
representação gráfica. Prestar assistência técnica no estudo e
desenvolvimento de projetos e pesquisas tecnológicas, voltadas às
atividades da área. Orientar e coordenar a execução dos serviços de
manutenção de equipamentos e instalações de ambientes e mobiliários
fixos. Reformar ambientes sem alteração estrutural.

Para atuação como Técnico em Design de Interiores, são fundamentais:

Conhecimentos interdisciplinares relacionados aos processos de criação,
envolvendo pesquisa, idealização, planejamento, execução técnica,
fruição e recepção estética. Competências comunicativas e empreendedoras
voltadas à proposição de projetos, ao coletivo, à gestão, à solução de
problemas e à resiliência, entre outras competências socioemocionais.

\textbf{CBOs não correspondidos e seus perfis:}

\textbf{CBO: 318005  ( Desenhista técnico )}

Analisa solicitação de desenho técnico, interpretando documentos de
apoio, tais como plantas, projetos, catálogos, croquis e normas. Propõe
ao solicitante alternativas para execução do desenho, chegando a um
acordo sobre os detalhes técnicos finais. Estima tempo para realização
do desenho e define prioridades, conforme cronograma. Programa ações
para realização do desenho técnico, observando características técnicas,
fazendo esboços, estabelecendo formatos e definindo escalas e sistemas
de representação. Elabora desenhos técnicos, representando formas
geométricas em duas e três dimensões. Utiliza computador, programa --
tal como CAD-Desenho Assistido por Computador (Computer Aided Design) -,
impressoras, plotters e outros periféricos, além de fazer uso de
instrumentos para execução do traçado diretamente em substratos de
papel. Constrói o desenho, traçando linhas auxiliares, cotando e
definindo a lista de materiais e componentes. Envia o desenho para
revisão, recebendo a aprovação do solicitante. Realiza cópias de
segurança e disponibiliza o desenho técnico final e/ou suas revisões
para as áreas afins, que providenciarão a entrega ao solicitante.
Conserva o local de trabalho limpo e organizado. Controla graus de
luminosidade do ambiente e verifica as condições ergonômicas para uso de
monitor de vídeo, mouse, mesa digitalizadora, prancheta e teclado.
Mantém equipamentos e instrumentos de trabalho em plenas condições de
uso e funcionamento. Requisita serviço de manutenção de computador,
impressoras, plotters e outros periféricos, quando necessário. Organiza
e arquiva documentos físicos e eletrônicos, tais como comprovantes de
aprovação e cópias de segurança dos desenhos técnicos. Destina
adequadamente os resíduos gerados, fazendo reciclagem e realizando
descarte de acordo com as normas ambientais.

\begin{Form}
\textbf{Esta correspondência está adequada?} \\ 
\ChoiceMenu[radio,radiosymbol=\ding{52},name=cbovalid_ 318005 ]{}{CBO deve ficar=sim, \hspace{6em} CBO não fica aqui=nao}
\end{Form}

\textbf{CBO: 318010  ( Desenhista copista )}

Analisa solicitação de cópia de desenho técnico, observando as
características técnicas do desenho original e interpretando documentos
de apoio, tais como originais de plantas, projetos, catálogos, croquis e
normas. Propõe ao solicitante alternativas para execução da cópia,
chegando a acordo em relação aos detalhes técnicos finais. Estima tempo
para realização do trabalho. Elabora cópias de desenhos técnicos,
reproduzindo representações de formas geométricas em duas e três
dimensões. Utiliza computador, programa -- tal como CAD-Desenho
Assistido por Computador (Computer Aided Design) -, mesas
digitalizadoras, impressoras, plotters e outros periféricos, além de
fazer uso de instrumentos para execução do traçado diretamente em
substratos de papel. Constrói a cópia do desenho, traçando linhas
auxiliares e cotando. Envia a cópia para revisão, recebendo a aprovação
do solicitante. Realiza cópias de segurança e disponibiliza o desenho
final copiado e/ou suas revisões para as áreas afins, que providenciarão
a entrega ao solicitante. Conserva o local de trabalho limpo e
organizado. Controla graus de luminosidade do ambiente e verifica as
condições ergonômicas para uso de monitor de vídeo, mouse, mesa
digitalizadora, prancheta e teclado. Mantém equipamentos e instrumentos
de trabalho em plenas condições de uso e funcionamento. Requisita
serviço de manutenção de computador, impressoras, plotters e outros
periféricos, quando necessário. Organiza e arquiva documentos físicos e
eletrônicos, tais como comprovantes de aprovação e uma via de cada
desenho copiado. Destina adequadamente os resíduos gerados, fazendo
reciclagem e realizando descarte de acordo com as normas ambientais.

\begin{Form}
\textbf{Esta correspondência está adequada?} \\ 
\ChoiceMenu[radio,radiosymbol=\ding{52},name=cbovalid_ 318010 ]{}{CBO deve ficar=sim, \hspace{6em} CBO não fica aqui=nao}
\end{Form}

\newpage
\section{ 09014 -- Técnico em Design de Móveis }

\textbf{Perfil do Curso (CNCT)}

O Técnico em Design de Móveis será habilitado para:

Desenvolver esboços, perspectivas e desenhos normatizados de móveis.
Realizar estudos volumétricos e maquetes convencionais e eletrônicas.
Aplicar aspectos ergonômicos ao projeto. Pesquisar e definir materiais,
ferragens e acessórios. Elaborar documentação técnica normatizada.
Acompanhar a execução de protótipos ou peças-piloto. Aplicar os
conceitos de sustentabilidade ao desenvolvimento de móveis.

Para atuação como Técnico em Design de Móveis, são fundamentais:

Conhecimentos interdisciplinares relacionados aos processos de criação,
envolvendo pesquisa, idealização, planejamento, execução técnica,
fruição e recepção estética. Competências comunicativas e empreendedoras
voltadas à proposição de projetos, ao coletivo, à gestão, à solução de
problemas e à resiliência, entre outras competências socioemocionais.

\textbf{CBOs não correspondidos e seus perfis:}

\textbf{CBO: 771120  ( Tanoeiro )}

Planeja o trabalho, interpretando projetos, desenhos e especificações
para a confecção e a restauração de pipas e outros vasilhames de madeira
e derivados. Seleciona matéria-prima, insumos e componentes, ferramentas
manuais, ferramentas elétricas portáteis, instrumentos de medição e
controle, máquinas e equipamentos. Define roteiro para as atividades de
confecção e restauração. Pode elaborar orçamento. Elabora esboço e
realiza dimensionamento do produto sob medida, conforme solicitação do
cliente e local de instalação. Confecciona gabaritos ou moldes para
execução das peças. Especifica madeiras, acessórios, ferragens e
acabamentos. Prepara máquinas e ferramentas. Confere parâmetros
operacionais e do produto, tais como tipo da madeira e seu nível de
umidade, posicionamento e geometria de corte, flexibilidade da madeira
para executar a curvatura das aduelas, tipo de chapa para fazer a cinta
do barril, entre outros. Confecciona pipas e outros vasilhames de
madeira e derivados, regulando as máquinas para obter o produto conforme
o projeto. Executa traçado em madeira, derivados e outros materiais,
observando o sentido dos veios, serrando, moldando e aplainando aduelas,
tampos e demais elementos, de acordo com as dimensões, curvaturas e
formas gerais requeridas. Monta pipas e outros vasilhames de madeira e
derivados, ajustando aduelas e tampos; utilizando elementos de união e
fixação das aduelas - tais como arcos ou argolas de metal provisórios e
definitivos -; aplicando massa para montagem, sob pressão; e submetendo
o conjunto a chamas, banhos de vapor ou outras formas de calor, para
dotar a madeira de mais flexibilidade. Realiza operações de acabamento,
apondo apliques, lâminas e anéis metálicos; colocando ferragens para
reajuste; regulando o funcionamento das partes móveis; e aplicando
produtos para correções, de acordo com o projeto. Realiza revestimento
interno no produto, para preservar as propriedades das substâncias
envasadas em contato com a madeira. Pode realizar teste de viscosidade
nos produtos utilizados para acabamento da superfície. Verifica os
produtos confeccionados, para garantir a qualidade de pipas e outros
vasilhames de madeira e derivados. Avalia matéria-prima, condições de
impermeabilidade, resistência, ausência de vazamentos, formas e
dimensões, encaixes, união e fixação das partes, acabamento das
superfícies internas e externas do recipiente e aplicação de acessórios.
Compara as características do produto final com os requisitos do cliente
ou do projeto, providenciando ajustes, quando necessários. Restaura
pipas e outros vasilhames de madeira e derivados, avaliando produtos
para determinar a viabilidade de restauração; confeccionando peças para
reposição; substituindo peças como aduelas, arcos metálicos e tampos
danificados; reapertando elementos de fixação; e preparando o produto
para acabamento. Prepara a expedição de pipas e outros vasilhames de
madeira e derivados confeccionados sob medida ou restaurados,
acondicionando-os para o transporte seguro e apondo instruções e
procedimentos de conservação do produto. Organiza o ambiente de trabalho
e os locais para carga, descarga e armazenamento de materiais. Realiza
limpeza e conservação dos recursos materiais. Mantém as máquinas, os
equipamentos e as ferramentas em plenas condições de uso e
funcionamento. Mantém os materiais, instrumentos e acessórios de
trabalho organizados e acondicionados. Controla desperdícios de
material. Separa, identifica e classifica resíduos -- sólidos e líquidos
-, para descarte, reúso e reciclagem. Utiliza equipamentos de proteção
individual e coletiva.

\begin{Form}
\textbf{Esta correspondência está adequada?} \\ 
\ChoiceMenu[radio,radiosymbol=\ding{52},name=cbovalid_ 771120 ]{}{CBO deve ficar=sim, \hspace{6em} CBO não fica aqui=nao}
\end{Form}

\newpage
\section{ 09019 -- Técnico em Figurino Cênico }

\textbf{Perfil do Curso (CNCT)}

O Técnico em Figurino Cênico será habilitado para:

Aplicar técnicas de pesquisa, concepção, desenho e execução de
figurinos. Aplicar técnicas de costura e modelagem de roupas. Produzir
figurinos (trajes e acessórios) de acordo com a época e o tema a ser
representado. Criar figurinos para personagens das artes cênicas,
audiovisual, dança e festas populares. Acompanhar as tendências
contemporâneas ligadas à criação de figurinos e as inovações
tecnológicas.

Para atuação como Técnico em Figurino Cênico, são fundamentais:

Conhecimentos interdisciplinares relacionados aos processos de criação,
envolvendo pesquisa, idealização, planejamento, execução técnica,
fruição e recepção estética. Competências comunicativas e empreendedoras
voltadas à proposição de projetos, ao coletivo, à gestão, à solução de
problemas e à resiliência, entre outras competências socioemocionais.

\textbf{CBOs não correspondidos e seus perfis:}

\textbf{CBO: 318810  ( Modelista de roupas )}

Planeja a confecção de moldes e o desenvolvimento de protótipos de
roupas, examinando a coleção do estilista; verificando tendências de
mercado; consultando boletins técnicos; avaliando estilos de design,
concepções e produtos de concorrentes; e coletando informações
tecnológicas em feiras e em desfiles de vestuário. Confecciona moldes de
roupas, interpretando modelos e desenhos; elaborando desenhos
planificados de roupas; traçando, cortando e codificando moldes. Utiliza
prancheta e softwares, tais como CAD-Desenho Assistido por Computador
(Computer Aided Design) e CAM--Manufatura Auxiliada por Computador
(Computer Aided Manufacturing). Marca referências em moldes e elabora
tabelas de medidas de roupas, para orientar a produção do vestuário.
Desenvolve protótipo de roupa, interpretando normas técnicas para
produção, elaborando ficha técnica, especificando aviamentos e
acessórios, ajustando moldes e gabaritos e determinando a posição de
ornamentos, detalhes e acessórios. Confecciona e testa peça-piloto
(protótipo). Requisita testes laboratoriais do protótipo. Realiza
ampliações e reduções das dimensões do produto por escala. Avalia
materiais para aquisição com vistas à produção de vestuário, verificando
a resistência dos materiais e analisando sua composição e acabamento.
Avalia a costurabilidade de tecidos e couros. Avalia a relação entre
custo e benefício de materiais. Interpreta resultados de testes
laboratoriais de materiais e elabora relatório.

\begin{Form}
\textbf{Esta correspondência está adequada?} \\ 
\ChoiceMenu[radio,radiosymbol=\ding{52},name=cbovalid_ 318810 ]{}{CBO deve ficar=sim, \hspace{6em} CBO não fica aqui=nao}
\end{Form}

\newpage
\section{ 09023 -- Técnico em Multimídia }

\textbf{Perfil do Curso (CNCT)}

O Técnico em Multimídia será habilitado para:

Desenvolver comunicação visual em meios eletrônicos, interfaces
interativas, publicações digitais, animações 2D e 3D, jogos eletrônicos,
web sites, web TV, TV digital e conteúdo audiovisual. Organizar e
preparar arquivos digitais para aplicações web e multimídia, animações e
games. Aplicar técnicas de tratamento de imagens estáticas e em
movimento que compõem estruturas de navegação em mídias digitais.
Executar atualização de páginas web e portais.

Para atuação como Técnico em Multimídia, são fundamentais:

Conhecimentos interdisciplinares relacionados aos processos de criação,
envolvendo pesquisa, idealização, planejamento, execução técnica,
fruição e recepção estética. Competências comunicativas e empreendedoras
voltadas à proposição de projetos, ao coletivo, à gestão, à solução de
problemas e à resiliência, entre outras competências socioemocionais.

\textbf{CBOs não correspondidos e seus perfis:}

\textbf{CBO: 317110  ( Desenvolvedor de sistemas de tecnologia da informação (técnico) )}

Desenvolve sistemas e aplicações de tecnologia da informação, elaborando
a interface gráfica e aplicando critérios ergonômicos de navegação.
Monta a estrutura do banco de dados e codifica os programas e
aplicativos, aplicando sistemas de rotinas de segurança. Compila os
programas. Testa aplicativos e programas, elaborando casos de testes.
Gera aplicativos para instalação e gerenciamento dos sistemas. Documenta
os sistemas e aplicações e avalia o desempenho dos produtos
desenvolvidos. Implanta sistemas e aplicações de tecnologia da
informação, instalando os programas e configurando os equipamentos em
que a aplicação será executada. Homologa os sistemas e aplicações
desenvolvidos, avaliando e validando os resultados da implantação.
Elabora material para capacitação de usuários. Avalia os objetivos e as
metas de projetos de sistemas e aplicações, tendo em vista os resultados
obtidos na fase de implantação. Publica o código final no servidor de
sistemas e aplicações. Realiza a manutenção de sistemas e aplicações de
tecnologia da informação, atualizando informações gráficas, textuais e
audiovisuais e convertendo seus códigos fontes para outras linguagens ou
plataformas. Implementa alterações nos sistemas e aplicações e em suas
estruturas de armazenamento de dados, adequando-os aos sistemas e
ambientes operacionais. Monitora a performance dos sistemas e
aplicações, utilizando ferramentas de software, para obter subsídios à
atividade de manutenção. Atualiza a documentação dos sistemas e
aplicações, em função das intervenções de manutenção. Fornece suporte
técnico para o cliente interno. Colabora na seleção dos recursos de
desenvolvimento, especificando máquinas, ferramentas, acessórios e
suprimentos. Participa da seleção das linguagens de programação e
metodologias. Seleciona ferramentas de desenvolvimento. Participa da
definição da nomenclatura padrão. Colabora no planejamento das etapas e
ações de trabalho, participando das definições de atividades, tarefas e
cronograma. Participa de reuniões com a equipe de trabalho ou com o
cliente. Participa das definições de padronizações. Projeta sistemas e
aplicações de tecnologia da informação, identificando a demanda do
cliente, coletando dados e desenvolvendo o leiaute de telas e
relatórios. Elabora o pré-projeto e os projetos conceitual, lógico,
estrutural, físico e gráfico dos sistemas e aplicações, definindo os
requisitos de projeto. Participa das definições dos critérios
ergonômicos de navegação e da interface de comunicação e interatividade.
Dimensiona a vida útil dos sistemas e aplicações, modela sua estrutura
de banco de dados e dimensiona a capacidade de armazenamento de dados
dos sistemas.

\begin{Form}
\textbf{Esta correspondência está adequada?} \\ 
\ChoiceMenu[radio,radiosymbol=\ding{52},name=cbovalid_ 317110 ]{}{CBO deve ficar=sim, \hspace{6em} CBO não fica aqui=nao}
\end{Form}

\newpage
\section{ 09030 -- Técnico em Produção de Áudio e Vídeo }

\textbf{Perfil do Curso (CNCT)}

O Técnico em Produção de Áudio e Vídeo será habilitado para:

Captar imagens e sons. Realizar ambientação e operação de equipamentos
por intermédio de recursos e linguagens. Investigar a utilização de
tecnologias de tratamento acústico, de imagem, luminosidade e animação.
Preparar material audiovisual. Elaborar fichas técnicas, mapas de
programação, distribuição, veiculação de produtos e serviços de
comunicação.

Para atuação como Técnico em Produção de Áudio e Vídeo, são
fundamentais:

Conhecimentos interdisciplinares relacionados aos processos de criação,
envolvendo pesquisa, idealização, planejamento, execução técnica,
fruição e recepção estética. Competências comunicacionais e
empreendedoras voltadas à proposição de projetos, ao coletivo, à gestão,
à solução de problemas e à resiliência, entre outras competências
socioemocionais.

\textbf{CBOs não correspondidos e seus perfis:}

\textbf{CBO: 374305  ( Operador de projetor cinematográfico )}

Planeja a operação de projetor cinematográfico, conferindo a programação
de cinema ou de outro local de exibição de filmes. Testa o funcionamento
do projetor. Realiza manutenção de primeiro nível, fazendo limpeza e, se
for o caso, lubrificação de equipamentos da cabine de projeção. Em
sistema analógico, verifica tipo de som e película e faz a montagem do
filme. Confere o nome e as partes do filme. Marca as partes do filme.
Monta o filme e os complementos. Revisa o filme e preenche o boletim de
revisão. Faz a seção teste, conferindo o tempo do filme e dos
complementos. Prepara a projeção de filme, em sistema analógico,
adaptando a lente e a máscara ao filme e o sistema de som. Instala o
filme no projetor. Inicia a projeção do filme, acionando os comandos do
equipamento. Configura e ajusta o projetor e a tela, para obter o
tamanho, a iluminação e o foco adequados das imagens. Regula o som, para
ter o volume e o tom adequados. Monitora a projeção e o volume de som.
Pode operar vários equipamentos e monitorar a projeção de uma película
em várias salas simultaneamente. Encerra a projeção do filme. Desmonta o
filme, separando suas partes e complementos. Rebobina o filme e cola
suas extremidades. Realiza o acondicionamento do filme em lata. Em
sistema digital, prepara a cabine de projeção, verificando itens do
ambiente -- como ventilação -- para instalação do projetor. Identifica
os conteúdos a serem projetados e cria a lista de reprodução
(``playlist'') da sala. Pode fazer a recepção do filme, transmitido via
internet e satélite ou recebido em disco rígido, para ser armazenado em
um servidor do cinema. Insere o arquivo digital HD-Alta Definição (High
Definition) - com imagem e som do filme - em servidor de vídeo,
localizado no projetor ou em aparelho conectado ao projetor por cabos de
HDMI-Interface Multimídia de Alta Definição (High-Definition Multimedia
Interface). Controla início e configuração do filme em tela de
computador ou outro dispositivo plugado ao projetor digital. Monitora a
projeção, tomando providências - como a troca de lâmpada do projetor
durante a operação -, para evitar eventuais falhas. Acompanha as
inovações na área, verificando o surgimento de novos equipamentos e
novas tecnologias digitais de projeção. Monitora a sala de cinema ou
outro ambiente de exibição de filmes, controlando a iluminação e o som.
Verifica o funcionamento do sistema de ar-condicionado. Detecta
problemas, corrigindo pequenas falhas ou providenciando serviço de
manutenção. Pode operar projetor de vídeo e sistema composto por dois ou
mais monitores que formam uma grande tela de exibição de vídeo
(``videowall''). Pode modificar o sistema de som para o vídeo. Conserva
a cabine de projeção limpa e organizada.

\begin{Form}
\textbf{Esta correspondência está adequada?} \\ 
\ChoiceMenu[radio,radiosymbol=\ding{52},name=cbovalid_ 374305 ]{}{CBO deve ficar=sim, \hspace{6em} CBO não fica aqui=nao}
\end{Form}

\textbf{CBO: 517330  ( Vigilante )}

Presta serviços de vigilância patrimonial - com a finalidade de garantir
a incolumidade física das pessoas e a integridade do patrimônio - a
instituições financeiras; a estabelecimentos comerciais, industriais e
de prestação de serviços; a residências; a entidades sem fins
lucrativos; e a órgãos e empresas públicas. Presta serviço de transporte
de numerário, bens ou valores, mediante a utilização de veículos, comuns
ou especiais. Realiza escolta armada de transporte de qualquer tipo de
carga ou de valor. Avalia condições da área de destino para desviar de
obstáculos e de obstruções no trajeto, percorrendo o que foi
preestabelecido. Controla objetos, cargas e veículos, verificando a
documentação da carga, acompanhando sua entrega no local de trabalho e
apreendendo objetos ilícitos ou irregulares. Identifica objetos achados
e perdidos para devolução. Presta serviço de segurança pessoal, com a
finalidade de garantir a incolumidade física de pessoas. Pode prevenir e
combater princípios de incêndio e prestar primeiros socorros. Faz uso de
equipamentos com tecnologias de segurança e monitoramento virtual. Pode
comunicar-se com a equipe por meio de gestos, sinais e códigos.
Comunica-se com a base durante a escolta. Relata ocorrências e interage
com órgãos oficiais.

\begin{Form}
\textbf{Esta correspondência está adequada?} \\ 
\ChoiceMenu[radio,radiosymbol=\ding{52},name=cbovalid_ 517330 ]{}{CBO deve ficar=sim, \hspace{6em} CBO não fica aqui=nao}
\end{Form}

\newpage
\section{ 09031 -- Técnico em Produção de Moda }

\textbf{Perfil do Curso (CNCT)}

O Técnico em Produção de Moda será habilitado para:

Coordenar a montagem de ambientes para divulgação da moda. Estabelecer
relação direta entre produto e consumidor por intermédio de catálogos,
desfiles, vídeos, fotografias e meios de comunicação em geral. Pesquisar
tendências de moda, de mercado e de lançamentos para construção de
estilos e composição visual. Elaborar a composição de looks para
apresentação pública de estilo, produção publicitária, vitrines,
exposições, desfiles.

Para atuação como Técnico em Produção de Moda, são fundamentais:

Conhecimentos interdisciplinares relacionados aos processos de criação,
envolvendo pesquisa, idealização, planejamento, execução técnica,
fruição e recepção estética. Competências comunicacionais e
empreendedoras voltadas à proposição de projetos, ao coletivo, à gestão,
à solução de problemas e à resiliência, entre outras competências
socioemocionais.

\textbf{CBOs não correspondidos e seus perfis:}

\textbf{CBO: 375110  ( Designer de vitrines )}

Realiza estudo preliminar, levantando as expectativas do cliente, os
limites orçamentários e o prazo para execução do projeto. Analisa o
perfil do público consumidor e as características do produto a ser
exposto. Realiza levantamento fotográfico e calcula as medidas do espaço
disponível, para identificar os elementos básicos para a concepção do
projeto da vitrine. Analisa a possibilidade de aproveitamento das
instalações pré-existentes. Concebe o projeto da vitrine, com base no
limite de custo e no prazo definidos. Desenha o croqui da vitrine,
utilizando técnicas de desenho de observação e perspectiva. Apresenta a
versão preliminar do projeto de vitrine ao cliente, podendo realizar
ajustes por ele solicitados. Elabora o projeto executivo, selecionando
elementos decorativos, cenográficos, mobiliários, suportes e
expositores. Escolhe a escala cromática. Aloca pontos elétricos e pontos
de iluminação. Representa graficamente as soluções para a vitrine,
utilizando desenho manual ou recursos tecnológicos, como CAD-Projeto
Assistido por Computador (Computer Aided Design). Faz o orçamento para
execução do projeto, a ser submetido à aprovação do cliente. Após
aprovação, consulta carteira de fornecedores, para aquisição de produtos
e serviços. Realiza a montagem da vitrine, de acordo com os princípios e
as técnicas de composição e exposição de produtos. Utiliza recursos de
iluminação, considerando os efeitos de destaque da decoração, dos
materiais promocionais e dos produtos expostos na vitrine. Avalia os
resultados do projeto, analisando atratividade da vitrine junto ao
público, variação de vendas do produto exposto e satisfação do cliente.
Realiza pesquisas nos mercados nacional e internacional e participa de
exposições, mostras e feiras, para conhecer inovações em vitrinismo.
Zela pela segurança, prevenindo acidentes e utilizando equipamentos de
proteção individual.

\begin{Form}
\textbf{Esta correspondência está adequada?} \\ 
\ChoiceMenu[radio,radiosymbol=\ding{52},name=cbovalid_ 375110 ]{}{CBO deve ficar=sim, \hspace{6em} CBO não fica aqui=nao}
\end{Form}

\newpage
\section{ 09033 -- Técnico em Rádio e Televisão }

\textbf{Perfil do Curso (CNCT)}

O Técnico em Rádio e Televisão será habilitado para:

Executar a produção e veiculação de programas radiofônicos e
televisivos. Realizar seleção musical, montagem de filmes, videotapes,
trilhas, vinhetas, jingles, spots, roteiros, locução e aplicação de
efeitos especiais. Operar equipamentos analógicos e digitais de estúdio
de gravação.

Para atuação como Técnico em Rádio e Televisão, são fundamentais:

Conhecimentos interdisciplinares relacionados aos processos de criação,
envolvendo pesquisa, idealização, planejamento, execução técnica,
fruição e recepção estética. Competências comunicativas e empreendedoras
voltadas à proposição de projetos, ao coletivo, à gestão, à solução de
problemas e à resiliência, entre outras competências socioemocionais.

\textbf{CBOs não correspondidos e seus perfis:}

\textbf{CBO: 374110  ( Técnico em instalação de equipamentos de áudio )}

Planeja a instalação de equipamentos de áudio, analisando projeto;
verificando as dimensões físicas do local do evento; examinando as
condições de infraestrutura para acesso e instalação de equipamentos; e
identificando os parâmetros acústicos do local. Estabelece as
necessidades técnicas, a partir das informações levantadas sobre o
evento de sonorização e gravação. Define ângulos de cobertura, níveis de
pressão sonora e resposta de frequência. Faz a especificação dos
equipamentos, tais como microfones, transdutores, cabos e conectores.
Elabora documento -- rider técnico -, relacionando as características
técnicas dos equipamentos de sonorização. Gera desenhos e diagramas de
instalação, incluindo atualizações referentes às alterações realizadas.
Orça o serviço de instalação de equipamentos para evento. Acondiciona
equipamentos, para transportá-los ao local do evento. Supervisiona o
transporte, o descarregamento e a limpeza dos equipamentos e acessórios.
Monta e testa os equipamentos de áudio e acessórios - tais como mesas
e/ou console de som analógico ou digital, caixas de som, monitores de
áudio, amplificadores, entre outros -, substituindo os que apresentam
defeitos. Executa o projeto de instalação de equipamentos, conferindo
sistema de energia elétrica e aterramento. Conecta sistema de caixas
acústicas, instala transdutores e conecta sinais de áudio. Confere o
funcionamento de sistema e a uniformidade de distribuição de áudio.
Checa os níveis de interferência em sistema. Configura e alinha sistema
de sonorização. Avalia as características - timbre, formato e tipo - de
fonte sonora. Seleciona e posiciona transdutores eletroacústicos. No
sistema de gravação, ajusta as estruturas de ganho de sistema e prepara
o sistema de monitoramento para gravação. Pode instalar, testar e
configurar equipamentos e softwares de captação e gravação de conteúdo
audiofônico para diversas mídias, considerando especificidades ligadas
aos canais de transmissão e distribuição -- radiodifusão, cabo, satélite
e rede mundial de computadores -- e respeitando o projeto previamente
estabelecido. Desconecta sistema de caixas acústicas, desinstala
transdutores e desconecta sinais de áudio, no final do evento. Prepara
os equipamentos para retirada do local do evento e para transporte.
Documenta as rotinas de trabalho, registrando ocorrências da produção,
relatando medidas de contingenciamento e comunicando ajustes aos
diferentes setores envolvidos nos processos de trabalho. Informa aos
profissionais, envolvidos em sonorização e gravação de áudio, sobre os
processos realizados de instalação, teste e configuração de
equipamentos. Auxilia na elaboração de pareceres e relatórios técnicos
de qualidade e segurança. Realiza manutenção de primeiro nível em
equipamentos e acessórios, executando pequenos reparos. Providencia
serviço de manutenção de equipamentos, para solucionar problemas mais
complexos. Trabalha com segurança, realizando o dimensionamento adequado
da alimentação elétrica; providenciando a fixação segura dos
equipamentos em paredes, tetos, treliças, pórticos e outros suportes;
examinando a localização e o posicionamento de equipamentos, de modo a
não obstruir saídas e passagens de emergência; verificando a
disponibilidade de dispositivos de extinção de incêndio adequados aos
equipamentos utilizados no sistema de áudio; e prevenindo acidentes, de
modo geral.

\begin{Form}
\textbf{Esta correspondência está adequada?} \\ 
\ChoiceMenu[radio,radiosymbol=\ding{52},name=cbovalid_ 374110 ]{}{CBO deve ficar=sim, \hspace{6em} CBO não fica aqui=nao}
\end{Form}

\textbf{CBO: 374115  ( Técnico em masterização de áudio )}

Planeja o serviço de masterização, programando as etapas a serem
realizadas e estabelecendo cronograma. Seleciona equipamentos analógicos
ou digitais. Faz o orçamento do serviço. Recebe arquivos de áudio após o
trabalho de mixagem, na forma de um arquivo único -- ``mix down'' -- ou
em ``stems'' separados. Analisa registros sonoros mixados, avaliando a
qualidade do produto. Realiza a restauração do áudio, marcando pontos
específicos e consertando pequenos erros, como chiados. Executa o
aprimoramento estéreo, acertando o equilíbrio espacial do áudio. Faz a
equalização, equilibrando e balanceando todos os elementos de áudio,
para soar bem em qualquer sistema de som. Melhora os elementos que
precisam se destacar e causar maior impacto em determinado momento do
produto sonoro. Efetua a compressão do áudio, uniformizando todas as
faixas. Passa a faixa de áudio por um tipo específico de compressor --
limitador (``limiter'') -, criando um pico teto e adequando o áudio com
base nesse elemento, para intensificar a música ou outro produto sonoro.
Na masterização de um álbum, organiza a sequência e o espaçamento (pausa
ou silêncio entre as faixas). Finaliza o processo de masterização,
armazenando a matriz de todas as cópias a serem feitas do produto
sonoro, para divulgação ao público em diversos formatos -- tais como
vinil, CD, fita e serviços de streaming (Spotify, iTunes, SoundCloud,
entre outros) - com a qualidade necessária para o ouvinte. Documenta as
rotinas de trabalho, organizando informações, registrando ocorrências da
produção, relatando medidas de contingenciamento e comunicando ajustes
aos diferentes setores envolvidos nos processos de trabalho. Auxilia na
elaboração de pareceres e relatórios técnicos de qualidade e segurança.
Realiza manutenção de primeiro nível em equipamentos e acessórios,
executando pequenos reparos. Providencia serviço de manutenção de
equipamentos, para solucionar problemas mais complexos. Trabalha com
segurança, prevenindo acidentes.

\begin{Form}
\textbf{Esta correspondência está adequada?} \\ 
\ChoiceMenu[radio,radiosymbol=\ding{52},name=cbovalid_ 374115 ]{}{CBO deve ficar=sim, \hspace{6em} CBO não fica aqui=nao}
\end{Form}

\textbf{CBO: 374120  ( Projetista de som )}

Prepara-se para elaborar projeto para sonorização e gravação do áudio de
produção audiovisual e/ou artística - tal como teatral, cinematográfica,
televisiva, entre outras -, colhendo informações. Interage com
profissionais envolvidos nas áreas técnica e artística da produção -
como diretor, diretor de arte, roteirista e editor -, para conhecer
roteiro, cenários, personagens e outros detalhes. Com relação ao local
da apresentação e/ou filmagem da produção, verifica as dimensões
físicas; examina as condições de infraestrutura para acesso e instalação
de equipamentos; e identifica os parâmetros acústicos. Elabora projeto
de sonorização e de gravação de áudio, planejando a infraestrutura
física e tecnológica; dimensionando e selecionando equipamentos de áudio
e softwares; configurando sistemas de sonorização; e escolhendo os
locais para captação e gravação do som. Elabora orçamento para
implementação do projeto. Realiza a implementação do projeto, após a
aprovação da proposta. Confere a uniformidade de distribuição de áudio.
Avalia as características - timbre, formato e tipo - de fonte sonora e
seleciona os transdutores eletroacústicos. Posiciona os microfones e
analisa os sinais de áudio. Mistura sinais de fontes de áudio. Distribui
sinais de áudio para outros sistemas. Executa o processo de gravação,
selecionando o meio de registro (mídia). Alinha os sistemas de gravação.
Arquiva o meio de registro, em ambiente que atenda às especificações
técnicas predefinidas. Realiza o tratamento dos registros sonoros,
selecionando-os e equalizando-os. Controla a dinâmica de registros
sonoros. Marca pontos específicos de faixas, rolos e arquivos digitais.
Ordena faixas e rolos em sequência predeterminada. Finaliza o produto
sonoro, para diversas mídias e/ou plataformas. Avalia a qualidade do
produto sonoro. Verifica a satisfação, dos responsáveis pela produção
audiovisual e/ou artística, com os resultados. Documenta o serviço
realizado, organizando informações, registrando ocorrências e relatando
interação com os setores envolvidos nos processos. Elabora relatório
sobre a implementação do projeto e sobre os resultados obtidos. Realiza
a manutenção de primeiro nível em equipamentos e acessórios. Providencia
serviço de manutenção de equipamentos, para solucionar problemas mais
complexos. Monitora a aplicação de técnicas de segurança durante a
implementação do projeto, examinando as medidas adotadas para prevenção
de acidentes.

\begin{Form}
\textbf{Esta correspondência está adequada?} \\ 
\ChoiceMenu[radio,radiosymbol=\ding{52},name=cbovalid_ 374120 ]{}{CBO deve ficar=sim, \hspace{6em} CBO não fica aqui=nao}
\end{Form}

\textbf{CBO: 374135  ( Projetista de sistemas de áudio )}

Elabora projeto de sistemas de sonorização e de gravação, levantando
informações para definir necessidades técnicas. Verifica dimensões
físicas e identifica parâmetros acústicos do local. Examina as condições
de infraestrutura para acesso e instalação de equipamentos. Define
ângulos de cobertura, níveis de pressão sonora e resposta de frequência.
Faz a especificação de equipamentos, microfones, transdutores, cabos e
conectores. Define a demanda de energia elétrica dos sistemas. Elabora
documento -- rider técnico -, relacionando as características técnicas
dos equipamentos de sonorização. Gera desenhos e diagramas de
instalação, incluindo atualizações referentes às alterações realizadas.
Orça o serviço de elaboração e implementação do projeto. Acompanha a
implementação do projeto, conferindo a uniformidade de distribuição de
áudio após a instalação dos equipamentos e acessórios. Configura e
alinha os sistemas de sonorização. Avalia as características - timbre,
formato e tipo - de fonte sonora e posiciona os transdutores
eletroacústicos. Analisa os sinais de áudio e, se adequados, os
distribui para outros sistemas. Ajusta as estruturas de ganho de sistema
e define padrões de sincronismo, no sistema de gravação. Monitora
sistemas, equipamentos e softwares de sonorização, captação e gravação
de sinais sonoros para diversas mídias e/ou plataformas. Avalia o
desempenho dos sistemas, realizando ajustes, quando necessário.
Documenta o serviço realizado, organizando informações, registrando
ocorrências e relatando interação com os setores envolvidos nos
processos. Elabora relatório sobre a implementação do projeto e sobre os
resultados obtidos. Realiza manutenção de primeiro nível em equipamentos
e acessórios, executando pequenos reparos. Providencia serviço de
manutenção de equipamentos, para solucionar problemas mais complexos.
Trabalha com segurança, prevenindo acidentes.

\begin{Form}
\textbf{Esta correspondência está adequada?} \\ 
\ChoiceMenu[radio,radiosymbol=\ding{52},name=cbovalid_ 374135 ]{}{CBO deve ficar=sim, \hspace{6em} CBO não fica aqui=nao}
\end{Form}

\textbf{CBO: 374140  ( Microfonista )}

Prepara os equipamentos de captação de sinais sonoros -- tais como
microfones ``boom'', direcionais, de vara, sem fio, de lapela, entre
outros -, acondicionando-os para transporte ao local do evento ou da
produção audiovisual e artística. Levanta informações, junto ao
responsável técnico do evento ou da produção -- teatral, televisiva,
musical ou outra --, sobre roteiro preestabelecido, personagens e
apresentadores. Configura, monta, posiciona e testa os equipamentos de
captação de sinais sonoros. Coloca e posiciona microfones sem fio e de
lapela em apresentadores e atores, considerando figurinos e acessórios
que serão utilizados, e verificando o funcionamento dos sistemas de
comunicação sem fio. Monitora o estado geral das baterias nos
equipamentos de captação de sinais sonoros, inventariando informações
relativas a cargas e vida útil das baterias, caixas sem fio,
equipamentos de retorno sonoro, entre outros. Participa de ensaios e
programas-piloto em diferentes gêneros e formatos de eventos e produções
audiovisuais e artísticas, verificando deslocamentos e movimentos de
câmera, ações de cena e ordem de falas, assim como a posição ideal de
microfones, de acordo com a iluminação ou a luz distribuída. Acompanha a
gravação ou a apresentação de evento ou da produção audiovisual e
artística, garantindo o cumprimento das diretrizes e orientações
estipuladas em roteiro e/ou apontamentos do responsável técnico. Envia
sinal de áudio para a câmera, se necessário, sincronizando o
``timecode'' e entregando sinal de áudio para os diferentes postos, tais
como diretor, continuísta e assistente de vídeo. Monitora e verifica
todo o sistema de comunicação sem fio utilizado. Confere e controla a
qualidade técnica do som dos equipamentos de captação de sinais sonoros,
analógicos e/ou digitais. Identifica eventuais defeitos ou problemas
técnicos, fazendo o reparo ou solicitando - para casos mais complexos -
o serviço de profissionais de manutenção. Documenta o serviço realizado,
organizando informações, registrando ocorrências e relatando interação
com os setores envolvidos nos processos. Auxilia na elaboração de
pareceres e relatórios técnicos. Trabalha com segurança, prevenindo
acidentes.

\begin{Form}
\textbf{Esta correspondência está adequada?} \\ 
\ChoiceMenu[radio,radiosymbol=\ding{52},name=cbovalid_ 374140 ]{}{CBO deve ficar=sim, \hspace{6em} CBO não fica aqui=nao}
\end{Form}

\newpage
\section{ 10001 -- Técnico em Açúcar e Álcool }

\textbf{Perfil do Curso (CNCT)}

O Técnico em Açúcar e Álcool será habilitado para:

Controlar e supervisionar operações de processos tecnológicos da
produção de açúcar e álcool e subprodutos. Realizar análises
físico-químicas e microbiológicas de matérias-primas e produtos dos
processos de industrialização da cana-de-açúcar para o controle de
qualidade. Desenvolver produtos e processos em açúcar e álcool.
Colaborar com equipe multidisciplinar nas fases de colheita, transporte,
moagem, industrialização e distribuição dos produtos e subprodutos.
Operar etapas e movimentação de materiais e insumos relacionados à área.

Para atuação como Técnico em Açúcar e Álcool, são fundamentais:

Conhecimentos e saberes relacionados aos processos de planejamento e
operação das atribuições da área, de modo a assegurar a saúde e a
segurança dos trabalhadores e dos futuros usuários e operadores de
empresas em processos de transformação em açúcar e álcool. Conhecimentos
e saberes relacionados à sustentabilidade do processo produtivo, às
normas e relatórios técnicos, à legislação da área, às novas tecnologias
relacionadas à indústria 4.0, à liderança de equipes, à solução de
problemas técnicos e à gestão de conflitos.

\textbf{CBOs não correspondidos e seus perfis:}

\textbf{CBO: 311115  ( Técnico em curtimento )}

Opera e controla máquinas e equipamentos empregados nas operações e
processos de curtimento de pele, recurtimento, matização, pré-acabamento
e acabamento de couro, conforme manuais técnicos e procedimentos de
produção. Monitora o funcionamento de máquinas, equipamentos e
instrumentos de medição e controle. Mantém máquinas e equipamentos em
condições de uso, realizando regulagens, quando necessárias. Realiza
análises químicas e físicas, utilizando métodos analíticos e
instrumentais, seguindo procedimentos e normas técnicas e de qualidade e
observando critérios de confiabilidade metrológica. Coleta e prepara
amostras. Prepara reagentes. Registra os resultados, assumindo
responsabilidade técnica por eles. Supervisiona processos de produção,
definindo e coordenando equipes de trabalho, elaborando o cronograma de
produção, definindo o fluxo do processo, monitorando os parâmetros do
processo e emitindo relatórios de produção. Participa do desenvolvimento
de produtos e processos, testando insumos e matérias primas, elaborando
formulações para produção, testando novos produtos, definindo o processo
de fabricação de produtos, fazendo a adequação dos produtos às
necessidades dos clientes e validando o produto desenvolvido. Participa
da implantação de programas de qualidade, analisando indicadores,
aplicando ferramentas da qualidade e implementando ações preventivas e
corretivas. Participa de auditorias. Colabora na implantação ou
reestruturação das instalações industriais, definindo leiautes e
fluxogramas de produção, acompanhando a montagem e instalação de
equipamentos e testando máquinas e equipamentos. Presta assistência
técnica, identificando necessidades dos clientes, diagnosticando
problemas técnicos e propondo soluções e melhorias de produtos e
processos. Elabora documentação técnica, emitindo laudos e redigindo
procedimentos e relatórios. Conserva o local de trabalho limpo e
organizado. Realiza tratamento de resíduos sólidos e efluentes gerados
no curtimento e no beneficiamento, para aproveitamento econômico -- tal
como aplicação na produção de biogás e de fertilizante agrícola -- e
diminuição do impacto ambiental.

\begin{Form}
\textbf{Esta correspondência está adequada?} \\ 
\ChoiceMenu[radio,radiosymbol=\ding{52},name=cbovalid_ 311115 ]{}{CBO deve ficar=sim, \hspace{6em} CBO não fica aqui=nao}
\end{Form}

\textbf{CBO: 325305  ( Técnico em biotecnologia )}

Prepara ambiente de trabalho, mapeando áreas de coleta e experimento;
selecionando equipamentos e utensílios; e montando utensílios de
laboratórios. Confere protocolos e procedimentos. Providencia materiais
aplicados à biotecnologia, coletando materiais para pesquisa;
selecionando e colocando identificação nos materiais biológico e
químico; descontaminando e esterilizando materiais. Faz armazenamento e
descarte dos materiais aplicados à biotecnologia de acordo com normas
técnicas e de biossegurança. Pode prestar apoio técnico à criação de
animais, para sua utilização em experimentos de pesquisa científica.
Cuida dos animais e monitora sua reprodução, sob a supervisão de
profissional de nível superior. Mensura taxa de natalidade e mortalidade
dos animais. Pode abater os animais em condições que envolvam, de acordo
com as espécies, o mínimo possível de sofrimento físico ou mental.
Controla as condições do laboratório para realização de experimento.
Confere condições ambientais, examinando luz, temperatura, umidade e
radiação. Dispõe os equipamentos, verificando seu funcionamento, fazendo
sua calibragem e realizando sua adequação ao experimento. Opera
equipamentos de laboratórios. Pode executar manutenção (de até segundo
nível) de equipamentos. Prepara meios de cultura e soluções, calculando
e pesando reagentes; misturando substâncias para produção de inóculos;
diluindo e concentrando soluções. Desenvolve culturas ``in vivo'' e ``in
vitro'' e marcadores moleculares, para estudos da fisiologia e da
bioquímica de células e para detecção da presença de doenças e
acompanhamento de sua evolução. Cultiva tecidos animal e vegetal para
multiplicação celular. Cultiva e inocula microrganismos. Macera tecidos
animais e vegetais. Realiza extração, replicação, sequenciamento e
quantificação de DNA-Ácido desoxirribonucleico (Deoxyribonucleic acid).
Analisa substâncias e compostos biológicos, medindo volume,
concentração, potencial hidrogeniônico e condutividade de soluções; e
determinando umidade de substâncias e compostos biológicos. Realiza
análises químicas e microbiológicas, atuando na produção de
imunobiológicos (medicamentos, soros e vacinas) e nos processos
biotecnológicos de indústrias - química, de alimentos, de medicamentos,
de cosméticos e de produtos veterinários -, sob supervisão de
profissional de nível superior. Controla a qualidade de matérias-primas
e produtos utilizados. Mantém-se atualizado em relação às pesquisas e
inovações na área de biotecnologia. Acompanha a introdução da
automatização e do uso de robótica nos processos e o desenvolvimento de
novos materiais e produtos. Monitora a limpeza do local de trabalho.
Orienta a aplicação de técnicas de higienização. Aplica procedimentos
para cumprimento de normas ambientais, promovendo o tratamento
estabelecido e o destino adequado de materiais biológicos e de outros
resíduos descartados. Atua seguindo as normas de segurança e de
biossegurança, submetendo-se a exames periódicos; vacinando-se contra
agentes patogênicos; e usando equipamentos de proteção individual.
Participa em ações para prevenção de acidentes, colaborando na
eliminação de situações de risco. Atua no combate a pequenos focos de
fogo, manuseando extintores de incêndio. Aplica cuidados de primeiros
socorros.

\begin{Form}
\textbf{Esta correspondência está adequada?} \\ 
\ChoiceMenu[radio,radiosymbol=\ding{52},name=cbovalid_ 325305 ]{}{CBO deve ficar=sim, \hspace{6em} CBO não fica aqui=nao}
\end{Form}

\textbf{CBO: 325310  ( Técnico em imunobiológicos )}

Prepara ambiente de trabalho, mapeando áreas de coleta e experimento;
selecionando equipamentos e utensílios; e montando utensílios de
laboratórios. Confere protocolos e procedimentos. Providencia materiais
aplicados à biotecnologia, coletando materiais para pesquisa;
selecionando e colocando identificação nos materiais biológico e
químico; descontaminando e esterilizando materiais. Faz armazenamento e
descarte dos materiais aplicados à biotecnologia de acordo com normas
técnicas e de biossegurança. Pode prestar apoio técnico à criação de
animais, para sua utilização em experimentos de pesquisa científica.
Cuida dos animais e monitora sua reprodução, sob a supervisão de
profissional de nível superior. Mensura taxas de natalidade e
mortalidade dos animais. Pode abater os animais em condições que
envolvam, de acordo com as espécies, o mínimo possível de sofrimento
físico ou mental. Controla as condições do laboratório para realização
de experimento. Confere condições ambientais, examinando luz,
temperatura, umidade e radiação. Dispõe os equipamentos, verificando seu
funcionamento, fazendo sua calibragem e realizando sua adequação ao
experimento. Opera equipamentos de laboratórios. Pode executar
manutenção (de até segundo nível) em equipamentos. Prepara meios de
cultura e soluções, calculando e pesando reagentes; misturando
substâncias para produção de inóculos; diluindo e concentrando soluções.
Desenvolve culturas ``in vivo'' e ``in vitro'' e marcadores moleculares,
para estudos da fisiologia e da bioquímica de células e para detecção da
presença de doenças e acompanhamento de sua evolução. Cultiva tecidos
animal e vegetal para multiplicação celular. Cultiva e inocula
microrganismos. Macera tecidos animais e vegetais. Faz extração,
replicação, sequenciamento e quantificação de DNA-Ácido
desoxirribonucleico (Deoxyribonucleic acid). Realiza manipulação de
RNA-Ácido ribonucleico (Ribonucleic acid). Analisa substâncias e
compostos biológicos, medindo volume, concentração, potencial
hidrogeniônico e condutividade de soluções; e determinando umidade de
substâncias e compostos biológicos. Realiza análises químicas e
microbiológicas, atuando na produção de imunobiológicos (medicamentos,
soros, vacinas e kits de diagnóstico) e no estudo de doenças infecciosas
causadas por microrganismos - tais como bactérias e vírus -, sob
supervisão de profissional de nível superior. Auxilia em pesquisas de
melhoramento genético. Pode colaborar em atividades de perícia criminal
e investigação genética. Controla a qualidade de matérias-primas e
produtos utilizados. Mantém-se atualizado em relação às pesquisas e
inovações na área de imunobiologia. Colabora em estudos para implantação
de inovações, tais como os ``órgãos-em-chip'', capazes de servir como
``simuladores'' de diferentes órgãos do corpo humano, permitindo que
pesquisadores monitorem a jornada de um corpo externo - como bactéria e
droga farmacêutica - com maior segurança e controle, em substituição ao
uso de animais em experimentos. Monitora a limpeza do local de trabalho.
Orienta a aplicação de técnicas de higienização. Aplica procedimentos
para cumprimento de normas ambientais, promovendo o tratamento
estabelecido e o destino adequado de materiais biológicos e de outros
resíduos descartados. Atua seguindo as normas de segurança e de
biossegurança, submetendo-se a exames periódicos; vacinando-se contra
agentes patogênicos; e usando equipamentos de proteção individual.
Participa em ações para prevenção de acidentes, colaborando na
eliminação de situações de risco. Atua no combate a pequenos focos de
fogo, manuseando extintores de incêndio. Aplica cuidados de primeiros
socorros.

\begin{Form}
\textbf{Esta correspondência está adequada?} \\ 
\ChoiceMenu[radio,radiosymbol=\ding{52},name=cbovalid_ 325310 ]{}{CBO deve ficar=sim, \hspace{6em} CBO não fica aqui=nao}
\end{Form}

\newpage
\section{ 10006 -- Técnico em Celulose e Papel }

\textbf{Perfil do Curso (CNCT)}

O Técnico em Celulose e Papel será habilitado para:

Controlar processos de obtenção da celulose e de fabricação de papel.
Realizar ensaios e análises químicas, físicas e físico-químicas de
matérias-primas e produtos seguindo normas e procedimentos técnicos.
Planejar, executar e supervisionar os processos de secagem e corte na
produção de papel.

Para atuação como Técnico em Celulose e Papel, são fundamentais:

Conhecimentos e saberes relacionados aos processos de planejamento e
operação das atribuições da área, de modo a assegurar a saúde e a
segurança dos trabalhadores e dos futuros usuários e operadores de
empresas em processos de obtenção da celulose e de fabricação de papel.
Conhecimentos e saberes relacionados à sustentabilidade do processo
produtivo, às normas e relatórios técnicos, à legislação da área, às
novas tecnologias relacionadas à indústria 4.0, à liderança de equipes,
à solução de problemas técnicos e à gestão de conflitos.

\textbf{CBOs não correspondidos e seus perfis:}

\textbf{CBO: 311115  ( Técnico em curtimento )}

Opera e controla máquinas e equipamentos empregados nas operações e
processos de curtimento de pele, recurtimento, matização, pré-acabamento
e acabamento de couro, conforme manuais técnicos e procedimentos de
produção. Monitora o funcionamento de máquinas, equipamentos e
instrumentos de medição e controle. Mantém máquinas e equipamentos em
condições de uso, realizando regulagens, quando necessárias. Realiza
análises químicas e físicas, utilizando métodos analíticos e
instrumentais, seguindo procedimentos e normas técnicas e de qualidade e
observando critérios de confiabilidade metrológica. Coleta e prepara
amostras. Prepara reagentes. Registra os resultados, assumindo
responsabilidade técnica por eles. Supervisiona processos de produção,
definindo e coordenando equipes de trabalho, elaborando o cronograma de
produção, definindo o fluxo do processo, monitorando os parâmetros do
processo e emitindo relatórios de produção. Participa do desenvolvimento
de produtos e processos, testando insumos e matérias primas, elaborando
formulações para produção, testando novos produtos, definindo o processo
de fabricação de produtos, fazendo a adequação dos produtos às
necessidades dos clientes e validando o produto desenvolvido. Participa
da implantação de programas de qualidade, analisando indicadores,
aplicando ferramentas da qualidade e implementando ações preventivas e
corretivas. Participa de auditorias. Colabora na implantação ou
reestruturação das instalações industriais, definindo leiautes e
fluxogramas de produção, acompanhando a montagem e instalação de
equipamentos e testando máquinas e equipamentos. Presta assistência
técnica, identificando necessidades dos clientes, diagnosticando
problemas técnicos e propondo soluções e melhorias de produtos e
processos. Elabora documentação técnica, emitindo laudos e redigindo
procedimentos e relatórios. Conserva o local de trabalho limpo e
organizado. Realiza tratamento de resíduos sólidos e efluentes gerados
no curtimento e no beneficiamento, para aproveitamento econômico -- tal
como aplicação na produção de biogás e de fertilizante agrícola -- e
diminuição do impacto ambiental.

\begin{Form}
\textbf{Esta correspondência está adequada?} \\ 
\ChoiceMenu[radio,radiosymbol=\ding{52},name=cbovalid_ 311115 ]{}{CBO deve ficar=sim, \hspace{6em} CBO não fica aqui=nao}
\end{Form}

\textbf{CBO: 311710  ( Colorista têxtil )}

Recebe, inspeciona e testa corantes, pigmentos, ligantes, resinas e
demais materiais utilizados na estamparia e no tingimento de materiais
têxteis, para garantir a conformidade com as especificações. Desenvolve
cores em laboratórios, selecionando, pesando e misturando pigmentos,
corantes e materiais auxiliares para preparar tinturas têxteis, de
acordo com especificações. Calcula e prepara receitas de cores para
produção, a partir de formulações de cores preestabelecidas. Coleta
amostras de materiais ou produtos intermediários e acabados para
realizar testes laboratoriais das receitas. Testa a qualidade de
tinturas aplicadas em materiais têxteis, comparando a cor com os padrões
estabelecidos - seja visualmente ou usando equipamentos como
colorímetros ou espectrofotômetros -, ajustando e padronizando cores.
Desenvolve cartela de cores para produtos têxteis e de vestuários,
pesquisando tendências, selecionando e harmonizando cores. Identifica
cores com códigos, em conformidade com padrões internacionais. Pode
interagir com o designer têxtil, para decidir as cores a serem usadas
para cada padrão de tecido e a combinação de cores em uma coleção para
os segmentos de mercado masculino, feminino e infantil. Analisa as
tendências da indústria da moda. Controla a qualidade de processo de
tingimento e estamparia, examinando as características físicas e
químicas das tinturas e corantes. Pode operar e manter máquinas e
equipamentos utilizados no processo de tingimento e estamparia de
materiais têxteis. Registra os dados operacionais ou de produção em
formulários especificados. Conserva local de trabalho limpo e
organizado. Armazena corantes, pigmentos, ligantes, resinas e demais
materiais. Destina adequadamente os resíduos gerados, reciclando
materiais e realizando descarte de acordo com as normas ambientais. Zela
pela segurança, prevenindo acidentes, eliminando condições inseguras e
utilizando equipamentos de proteção individual.

\begin{Form}
\textbf{Esta correspondência está adequada?} \\ 
\ChoiceMenu[radio,radiosymbol=\ding{52},name=cbovalid_ 311710 ]{}{CBO deve ficar=sim, \hspace{6em} CBO não fica aqui=nao}
\end{Form}

\textbf{CBO: 311715  ( Preparador de tintas )}

Recebe, inspeciona e testa corantes, pigmentos e demais materiais
utilizados na preparação das tintas, para garantir a conformidade com as
especificações. Seleciona, pesa e mistura pigmentos, corantes e
materiais auxiliares para preparar tintas, de acordo com fórmula
preestabelecida. Faz a mistura manual ou mecanicamente. Coleta amostras
de materiais ou produtos intermediários e acabados para testes
laboratoriais. Testa a qualidade de tintas aplicadas em materiais ou
produtos, comparando a cor com os padrões estabelecidos - seja
visualmente ou usando equipamentos como colorímetros ou
espectrofotômetros -, ajustando e padronizando cores. Controla a
qualidade de processo de preparação de tintas. Examina as
características físicas e químicas das tintas, tais como viscosidade,
pH, densidade, aderência e espessura de camada de tintas aplicadas.
Opera e mantém máquinas utilizadas no processo de preparação de tintas.
Pode usar software específico para modelamento dos projetos de aplicação
das tintas. Pode preparar biotintas, tintas autorregenerativas, bem como
utilizar nanopigmentos, nanofibras e nanoativos. Presta assistência ao
cliente interno e externo. Registra os dados operacionais ou de produção
em formulários especificados. Conserva local de trabalho limpo e
organizado. Armazena corantes, pigmentos e demais materiais utilizados
na preparação das tintas. Destina adequadamente os resíduos gerados,
reciclando materiais e realizando descarte de acordo com as normas
ambientais. Zela pela segurança, prevenindo acidentes, eliminando
condições inseguras e utilizando equipamentos de proteção individual e
coletiva.

\begin{Form}
\textbf{Esta correspondência está adequada?} \\ 
\ChoiceMenu[radio,radiosymbol=\ding{52},name=cbovalid_ 311715 ]{}{CBO deve ficar=sim, \hspace{6em} CBO não fica aqui=nao}
\end{Form}

\newpage
\section{ 10008 -- Técnico em Construção Naval }

\textbf{Perfil do Curso (CNCT)}

O Técnico em Construção Naval será habilitado para:

Realizar ensaios e testes e montar componentes na fabricação e
manutenção naval. Desenvolver projetos de construção naval. Controlar e
inspecionar os processos de construção em plantas navais. Apoiar a
coordenação de construção de embarcações e estruturas hidroviárias.
Realizar manutenção e operação de sistemas de navegação. Selecionar
materiais a serem empregados na construção naval. Analisar custos
operacionais na área. Testar a velocidade e a segurança de barcos e
navios. Montar e organizar estaleiros. Operar sistemas de logística para
controle do frete, do armazenamento e da distribuição de cargas. Emitir
laudos técnicos e fazer vistorias nas companhias de navegação dentro de
suas atribuições técnicas.

Para atuação como Técnico em Construção Naval, são fundamentais:

Conhecimentos e saberes relacionados aos processos de planejamento e
operação das atribuições da área, de modo a assegurar a saúde e a
segurança dos trabalhadores e dos futuros usuários e operadores de
empresas em processos de transformação em construção naval.
Conhecimentos e saberes relacionados à sustentabilidade do processo
produtivo, às normas e relatórios técnicos, à legislação da área, às
novas tecnologias relacionadas à indústria 4.0, à liderança de equipes,
à solução de problemas técnicos e à gestão de conflitos.

\textbf{CBOs não correspondidos e seus perfis:}

\textbf{CBO: 318205  ( Desenhista técnico mecânico )}

Planeja o trabalho, levando em consideração o cronograma geral da área
ou do projeto -- ao qual se vinculam as atividades de elaboração do
desenho técnico mecânico --, estabelecendo prioridades e estimando o
tempo de execução do desenho. Define os meios computacionais e/ou
manuais a serem empregados nos trabalhos e elabora o cronograma
específico de execução do desenho. Coleta os dados necessários à
elaboração do desenho técnico mecânico, estudando as características do
projeto; identificando interfaces com áreas envolvidas no projeto; e
levantando informações básicas e informações complementares junto às
áreas de interface, a fim de atender aos requisitos de projeto em todas
as etapas de elaboração do desenho. Consulta tabelas, gráficos, normas
técnicas e outras fontes. Identifica o material -- das peças ou dos
conjuntos que integram os artefatos mecânicos a serem desenhados -- e
identifica a aplicabilidade do projeto. Identifica inovações
tecnológicas na área da mecânica e seus reflexos nos elementos de
composição do desenho. Elabora croquis, com o auxílio de instrumentos de
medição e materiais de desenho, para fazer estudo prévio do desenho
definitivo. Elabora desenhos técnicos detalhados de componentes,
ferramentas, dispositivos, máquinas, equipamentos, motores, veículos e
outros artefatos mecânicos, utilizando programas computacionais de apoio
e/ou instrumentos manuais e materiais de desenho. Utiliza programas do
tipo CAD-Desenho Assistido por Computador (Computer Aided Design) e os
recursos -- de hardware e de software -- associados a essas ferramentas.
Pode executar o desenho técnico mecânico a partir de informações da
representação de artefatos mecânicos em três dimensões (3D), utilizando
ferramentas computacionais integradas de representação em 3D e de
simulações. Na execução do desenho, aplica normas para desenho técnico e
realiza os cálculos necessários à execução do desenho, tais como cálculo
das dimensões, das superfícies, do volume, do peso e do centro de
gravidade de cada peça. Pode utilizar recursos de software para auxílio
ao cálculo. Efetua cálculos para estabelecer relações entre as
diferentes partes do artefato mecânico desenhado, dimensionando os
elementos parciais em escalas adequadas. Executa a representação de
todos os elementos do desenho com precisão técnica, clareza e nitidez.
Detalha desenho de conjunto e desenho de montagem. Elabora lista de
peças de conjunto. Confere o desenho e submete o esboço elaborado à
apreciação do autor do projeto, consultando-o sobre possíveis correções
ou alterações, para efetuar os ajustes necessários. Elabora o desenho
definitivo, registrando comentários, instruções gerais e informações
sobre materiais, processos e técnicas de fabricação e montagem. Faz
cópias do desenho, por meio de plotters, impressoras eletrônicas ou
máquinas copiadoras de diferentes tecnologias. Acompanha fabricação e
montagem dos artefatos mecânicos desenhados, assistindo o executante na
fabricação e conferindo as peças, os conjuntos e a montagem em relação
ao desenho. Pode detectar erros, propor modos de corrigi-los e sugerir
melhorias. Elabora relatórios de acompanhamento dos trabalhos. Pode
alterar desenhos, em função de mudanças de projeto, adequação às normas
ou outras demandas. Pode elaborar desenhos para manuais de referência e
publicações técnicas, a fim de descrever a operação e a manutenção de
artefatos mecânicos. Identifica os desenhos para arquivo. Organiza os
arquivos dos desenhos, em meios eletrônicos e em substratos de papel ou
similares. Controla a distribuição de cópias dos desenhos. Conserva o
local de trabalho limpo e organizado. Controla luminosidade do ambiente
e verifica as condições ergonômicas. Destina adequadamente os resíduos
gerados. Mantém equipamentos e instrumentos de trabalho em plenas
condições de uso e funcionamento. Zela pela segurança, atuando na
prevenção de acidentes e utilizando equipamentos de proteção individual.

\begin{Form}
\textbf{Esta correspondência está adequada?} \\ 
\ChoiceMenu[radio,radiosymbol=\ding{52},name=cbovalid_ 318205 ]{}{CBO deve ficar=sim, \hspace{6em} CBO não fica aqui=nao}
\end{Form}

\textbf{CBO: 318210  ( Desenhista técnico aeronáutico )}

Planeja o trabalho, levando em consideração o cronograma geral da área
ou do projeto -- ao qual se vinculam as atividades de elaboração do
desenho técnico aeronáutico --, estabelecendo as prioridades e estimando
o tempo de execução do desenho. Define os meios computacionais e/ou
manuais a serem empregados nos trabalhos e elabora o cronograma
específico de execução do desenho. Coleta os dados necessários à
elaboração do desenho técnico aeronáutico, estudando as características
do projeto; identificando interfaces com áreas envolvidas no projeto; e
levantando informações básicas e informações complementares junto às
áreas de interface, a fim de atender aos requisitos de projeto em todas
as etapas de elaboração do desenho. Consulta tabelas, gráficos, normas
técnicas e outras fontes. Identifica o material -- das peças ou dos
conjuntos que integram os artefatos mecânicos de aeronaves a serem
desenhados -- e identifica a aplicabilidade do projeto. Identifica
inovações tecnológicas na área aeronáutica e seus reflexos nos elementos
de composição do desenho. Elabora croquis, com o auxílio de instrumentos
de medição e materiais de desenho, para fazer estudo prévio do desenho
definitivo. Elabora desenhos técnicos detalhados de estruturas,
componentes, equipamentos, sistemas e outros artefatos mecânicos de
células de aeronaves, utilizando programas computacionais de apoio e/ou
instrumentos manuais e materiais de desenho. Utiliza programas do tipo
CAD-Desenho Assistido por Computador (Computer Aided Design) e os
recursos -- de hardware e de software -- associados a essas ferramentas.
Pode executar o desenho a partir de informações da representação de
artefatos mecânicos de aeronaves em três dimensões (3D), utilizando
ferramentas computacionais integradas de representação em 3D e de
simulações. Na execução do desenho, aplica normas para desenho técnico e
realiza os cálculos necessários à execução do desenho, tais como cálculo
das dimensões, das superfícies, do volume, do peso e do centro de
gravidade de cada peça. Pode utilizar recursos de software para auxílio
ao cálculo. Efetua cálculos para estabelecer relações entre as
diferentes partes do artefato mecânico de aeronave desenhado,
dimensionando os elementos parciais em escalas adequadas. Executa a
representação de todos os elementos do desenho com precisão técnica,
clareza e nitidez. Detalha desenho de conjunto e desenho de montagem.
Elabora lista de peças de conjunto. Confere o desenho e submete o esboço
elaborado à apreciação do autor do projeto, consultando-o sobre
possíveis correções ou alterações, para efetuar os ajustes necessários.
Elabora o desenho definitivo, registrando comentários, instruções gerais
e informações sobre materiais, processos e técnicas de fabricação e
montagem. Faz cópias do desenho, por meio de plotters, impressoras
eletrônicas ou máquinas copiadoras de diferentes tecnologias. Acompanha
fabricação e montagem dos artefatos mecânicos de aeronaves desenhados,
assistindo o executante na fabricação e conferindo as peças, os
conjuntos e a montagem em relação ao desenho. Pode detectar erros,
propor modos de corrigi-los e sugerir melhorias. Elabora relatórios de
acompanhamento dos trabalhos. Pode alterar desenhos, em função de
mudanças de projeto, adequação às normas ou outras demandas. Pode
elaborar desenhos para manuais de referência e publicações técnicas, a
fim de descrever a operação e a manutenção de artefatos mecânicos de
aeronaves. Identifica os desenhos para arquivo. Organiza os arquivos dos
desenhos em substratos de papel ou similares e em meios eletrônicos.
Controla a distribuição de cópias dos desenhos. Conserva o local de
trabalho limpo e organizado. Controla luminosidade do ambiente e
verifica as condições ergonômicas. Destina adequadamente os resíduos
gerados. Mantém equipamentos e instrumentos de trabalho em plenas
condições de uso e funcionamento. Zela pela segurança, atuando na
prevenção de acidentes e utilizando equipamentos de proteção individual.

\begin{Form}
\textbf{Esta correspondência está adequada?} \\ 
\ChoiceMenu[radio,radiosymbol=\ding{52},name=cbovalid_ 318210 ]{}{CBO deve ficar=sim, \hspace{6em} CBO não fica aqui=nao}
\end{Form}

\newpage
\section{ 10020 -- Técnico em Têxtil }

\textbf{Perfil do Curso (CNCT)}

O Técnico em Têxtil será habilitado para:

Supervisionar os processos produtivos na cadeia têxtil, da ?ação ao
bene?ciamento. Planejar e controlar as operações nos processos nas áreas
de ?ação, tecelagem e bene?ciamento têxtil. Desenvolver padronagens de
malharia ou tecido plano. Desenvolver produtos e processos de
tinturaria, estamparia e acabamento ?nal. Realizar testes de controle de
qualidade, químicos, físicos e colorimétricos. Analisar laudos técnicos.
Controlar estoques de produtos acabados.

Para atuação como Técnico em Têxtil, são fundamentais:

Conhecimentos e saberes relacionados aos processos de planejamento e
operação das atribuições da área, de modo a assegurar a saúde e a
segurança dos trabalhadores e dos futuros usuários e operadores de
empresas em processos de transformação têxtil. Conhecimentos e saberes
relacionados à sustentabilidade do processo produtivo, às normas e
relatórios técnicos, à legislação da área, às novas tecnologias
relacionadas à indústria 4.0, à liderança de equipes, à solução de
problemas técnicos e à gestão de conflitos.

\textbf{CBOs não correspondidos e seus perfis:}

\textbf{CBO: 311705  ( Colorista de papel )}

Recebe, inspeciona e testa corantes, pigmentos e demais materiais
utilizados na coloração de papéis, para garantir a conformidade com as
especificações. Desenvolve cores em laboratórios, selecionando, pesando
e misturando pigmentos, corantes e materiais auxiliares, para preparar
tinturas de acordo com especificações. Coleta amostras de materiais ou
produtos intermediários e acabados para testes laboratoriais. Testa a
qualidade de tinturas aplicadas em papel, comparando a cor com os
padrões estabelecidos - seja visualmente ou usando equipamentos como
colorímetros ou espectrofotômetros -, ajustando e padronizando cores.
Desenvolve cores e elabora cartelas de amostras de cores de acordo com
tendências e em conformidade com padrões internacionais. Controla a
qualidade do processo de tingimento e impressão, verificando a
viscosidade e o pH de tintas pastosas e examinando a densidade e a
espessura de camada de tinta nas impressões. Pode operar e manter
máquinas utilizadas no processo de coloração de papel. Registra os dados
operacionais ou de produção em formulários especificados. Conserva local
de trabalho limpo e organizado. Armazena corantes, pigmentos e demais
materiais. Destina adequadamente os resíduos gerados, reciclando papel e
realizando descarte de acordo com as normas ambientais. Zela pela
segurança, prevenindo acidentes, eliminando condições inseguras e
utilizando equipamentos de proteção individual.

\begin{Form}
\textbf{Esta correspondência está adequada?} \\ 
\ChoiceMenu[radio,radiosymbol=\ding{52},name=cbovalid_ 311705 ]{}{CBO deve ficar=sim, \hspace{6em} CBO não fica aqui=nao}
\end{Form}

\textbf{CBO: 311715  ( Preparador de tintas )}

Recebe, inspeciona e testa corantes, pigmentos e demais materiais
utilizados na preparação das tintas, para garantir a conformidade com as
especificações. Seleciona, pesa e mistura pigmentos, corantes e
materiais auxiliares para preparar tintas, de acordo com fórmula
preestabelecida. Faz a mistura manual ou mecanicamente. Coleta amostras
de materiais ou produtos intermediários e acabados para testes
laboratoriais. Testa a qualidade de tintas aplicadas em materiais ou
produtos, comparando a cor com os padrões estabelecidos - seja
visualmente ou usando equipamentos como colorímetros ou
espectrofotômetros -, ajustando e padronizando cores. Controla a
qualidade de processo de preparação de tintas. Examina as
características físicas e químicas das tintas, tais como viscosidade,
pH, densidade, aderência e espessura de camada de tintas aplicadas.
Opera e mantém máquinas utilizadas no processo de preparação de tintas.
Pode usar software específico para modelamento dos projetos de aplicação
das tintas. Pode preparar biotintas, tintas autorregenerativas, bem como
utilizar nanopigmentos, nanofibras e nanoativos. Presta assistência ao
cliente interno e externo. Registra os dados operacionais ou de produção
em formulários especificados. Conserva local de trabalho limpo e
organizado. Armazena corantes, pigmentos e demais materiais utilizados
na preparação das tintas. Destina adequadamente os resíduos gerados,
reciclando materiais e realizando descarte de acordo com as normas
ambientais. Zela pela segurança, prevenindo acidentes, eliminando
condições inseguras e utilizando equipamentos de proteção individual e
coletiva.

\begin{Form}
\textbf{Esta correspondência está adequada?} \\ 
\ChoiceMenu[radio,radiosymbol=\ding{52},name=cbovalid_ 311715 ]{}{CBO deve ficar=sim, \hspace{6em} CBO não fica aqui=nao}
\end{Form}

\textbf{CBO: 311725  ( Tingidor de couros e peles )}

Recebe, inspeciona e testa corantes, pigmentos e demais materiais
utilizados no tingimento de couros e peles, para garantir a conformidade
com as especificações. Seleciona, pesa e mistura pigmentos, corantes e
materiais auxiliares para preparar tinturas e banhos de acordo com
especificações. Calcula receitas para produção a partir de formulações
de cores preestabelecidas. Coleta amostras de materiais ou produtos
intermediários e acabados para testes laboratoriais. Testa a qualidade
de tinturas aplicadas em couros e peles, comparando a cor com os padrões
estabelecidos - seja visualmente ou usando equipamentos como
colorímetros ou espectrofotômetros -, ajustando e padronizando cores.
Controla a qualidade do processo de tingimento, verificando as
características físicas e químicas das tinturas e banhos. Opera e mantém
máquinas e equipamentos utilizados no processo de tingimento de couros e
peles. Registra os dados operacionais ou de produção em formulários
especificados. Conserva local de trabalho limpo e organizado. Armazena
corantes, pigmentos e demais materiais utilizados no tingimento de
couros e peles. Trata resíduos e efluentes gerados no processo de
tingimento, de acordo com as normas ambientais. Zela pela segurança,
prevenindo acidentes, eliminando condições inseguras e utilizando
equipamentos de proteção individual e coletiva.

\begin{Form}
\textbf{Esta correspondência está adequada?} \\ 
\ChoiceMenu[radio,radiosymbol=\ding{52},name=cbovalid_ 311725 ]{}{CBO deve ficar=sim, \hspace{6em} CBO não fica aqui=nao}
\end{Form}

\textbf{CBO: 721305  ( Afiador de cardas )}

Planeja o trabalho de afiação, consultando normas, instruções e
parâmetros. Estuda desenhos e especificações das cardas para determinar
os procedimentos de afiação, bem como planejar configurações da máquina
e sequências operacionais. Seleciona e monta rebolos ou lixas nas
máquinas, de acordo com as especificações, usando ferramentas manuais e
aplicando conhecimentos sobre abrasivos. Seleciona e extrai as pontas
desgastadas da guarnição da carda a serem afiadas, utilizando
ferramentas adequadas, ou prepara para a afiação diretamente na
guarnição, fazendo uso de máquinas especiais. Escolhe angulação, de
acordo com o material a ser afiado. Afia por processo manual ou
semiautomático, interpretando instruções de trabalho. Prepara a
ferramenta de afiação, executando o dressamento do rebolo, instalando-o
na máquina de afiar e fazendo os ajustes necessários. Instala lixas na
máquina de afiar. Verifica refrigeração e lubrificação do material a ser
afiado e regula a máquina para afiar. Gira as válvulas para direcionar o
fluxo de líquido de refrigeração durante o processo de afiação. Aciona a
máquina, manejando os controles de partida e de movimentação manual ou
semiautomática, para colocar o rebolo ou a lixa em contato com a peça,
procedendo à afiação de cada ponta. Controla o processo, avaliando e
compensando desgastes da ferramenta de afiação e a substituindo, se
necessário, quando desgastada. Pode trabalhar com ferramentas manuais.
Afia por processo automático, preparando, configurando e programando
máquinas de afiação de cardas convencionais e de alto desempenho. Coloca
as peças -- pontas ou guarnições completas das cardas -- a serem afiadas
na máquina. Aciona comandos para dar início às operações de afiação.
Monitora o processo. Verifica o acabamento final da afiação. Troca
instrumentos de afiação desgastados quando atingem condições mínimas de
corte. Retira a peça afiada, detendo a máquina e utilizando ferramentas
apropriadas. Reinstala a ponta afiada na guarnição da carda. Pode
retificar guarnição de cardas. Controla a qualidade, verificando as
dimensões das peças ou lotes no início do processo e definindo a
quantidade - totalidade ou amostras das peças - a inspecionar.
Identifica os pontos a inspecionar e realiza a inspeção do produto,
visualmente ou fazendo uso de equipamentos. Examina e mede as peças para
garantir que as superfícies e as dimensões atendam às especificações ao
final do processo. Identifica material a ser refugado ou retrabalhado.
Executa manutenção básica - como limpeza e lubrificação de peças - das
máquinas e ferramentas de afiação. Remove e substitui peças gastas ou
quebradas, usando ferramentas manuais. Trabalha com segurança, seguindo
normas de segurança e de medicina do trabalho. Ajusta as condições
ergonômicas. Utiliza equipamentos de proteção individual. Participa do
levantamento de riscos ambientais e de ações preventivas contra
incêndios e acidentes. Presta primeiros socorros. Conserva o ambiente de
trabalho limpo e organizado. Destina adequadamente os resíduos gerados,
reciclando materiais e realizando descarte de acordo com as normas
ambientais.

\begin{Form}
\textbf{Esta correspondência está adequada?} \\ 
\ChoiceMenu[radio,radiosymbol=\ding{52},name=cbovalid_ 721305 ]{}{CBO deve ficar=sim, \hspace{6em} CBO não fica aqui=nao}
\end{Form}

\newpage
\section{ 11016 -- Técnico em Recursos Pesqueiros }

\textbf{Perfil do Curso (CNCT)}

O Técnico em Recursos Pesqueiros será habilitado para:

Realizar operações do setor pesqueiro com base no manejo e na qualidade
da cadeia produtiva do pescado. Analisar e avaliar os aspectos técnicos,
sociais e econômicos da cadeia produtiva do setor pesqueiro. Prestar
assistência técnica e assessoria ao estudo e ao desenvolvimento de
projetos e pesquisas tecnológicas, ou aos trabalhos de vistoria,
perícia, arbitramento e consultoria. Elaborar orçamentos, laudos,
pareceres, relatórios e projetos, inclusive de incorporação de novas
tecnologias. Analisar a viabilidade técnica e econômica de propostas de
projetos aquícolas e pesqueiros. Prestar assistência técnica às áreas de
crédito rural, agroindustrial e impacto ambiental. Reconhecer os
aspectos fisiológicos, biológicos e ecológicos das espécies cultivadas e
exploradas pela pesca. Planejar e organizar ações de extensão pesqueira.
Planejar, organizar, dirigir e controlar as operações de pesca e de
despesca. Confeccionar, montar e operar apetrechos e equipamentos de
pesca e de aquicultura e dar manutenção adequada. Realizar procedimentos
para reprodução das principais espécies de interesse pesqueiro.
Monitorar e realizar o beneficiamento e o processamento de pescado em
embarcações pesqueiras e frigoríficos de pescado. Aplicar boas práticas
de manipulação e fabricação, supervisionar as etapas de conservação,
processamento, beneficiamento e comercialização do pescado. Elaborar,
aplicar e monitorar programas profiláticos, higiênicos e sanitários na
cadeia produtiva do pescado. Implantar e gerenciar sistemas de controle
de qualidade na produção do pescado. Emitir laudos e documentos de
classificação e exercer a fiscalização de produtos pesqueiros. Treinar e
conduzir equipes nas suas modalidades de atuação profissional. Prevenir
situações de risco à segurança no trabalho. Aplicar as legislações
pertinentes ao processo produtivo e ao meio ambiente. Identificar e
aplicar técnicas mercadológicas para distribuição e comercialização do
pescado. Executar a gestão econômica e financeira da produção pesqueira.
Administrar e gerenciar empreendimentos do setor pesqueiro. Conduzir
embarcações de pesca. Operar equipamentos como radares, bússolas,
barômetros e de Sistema de Navegação Global por Satélite (GNSS).

Para a atuação como Técnico em Recursos Pesqueiros, são fundamentais:

Conhecimentos e saberes relacionados a processos de gestão de negócios
voltados à aquicultura e pesca, à legislação pesqueira brasileira, à
legislação ambiental, ao planejamento de produção, à gestão de projetos,
à gestão de processos, ao empreendedorismo, ao mercado e comercialização
do pescado, à extensão aquícola e pesqueira, à aquicultura em águas da
União, a políticas públicas para o desenvolvimento da aquicultura e da
pesca, ao associativismo e ao cooperativismo. Domínio do uso de
tecnologias da informação.\\
Habilidade de comunicação, resolução de situações-problema, trabalho em
equipe e gestão de conflitos.

\textbf{CBOs não correspondidos e seus perfis:}

\textbf{CBO: 321310  ( Técnico em carcinicultura )}

Elabora projeto de implantação e de operação de sistema de criação de
camarões marinhos ou de camarões de água doce em cativeiro. Faz análise
da viabilidade técnica e econômica do projeto. Apresenta as etapas de
implantação. Relaciona máquinas, equipamentos, ferramentas, instrumentos
e materiais necessários. Detalha tamanho de equipe de trabalho e suas
atribuições. Elabora orçamento. Presta orientações técnicas na
implantação do projeto de carcinicultura, conforme características da
área, do clima e da localização geográfica. Seleciona sistema extensivo,
semi-intensivo ou intensivo de criação. Recomenda espécie de camarões
para criação. Orienta a construção e a adaptação de instalações - no mar
ou na terra - para a criação de camarões. Executa levantamento
topográfico. Indica necessidade de drenagem, colocação de telas e outras
providências. Orienta a montagem de sistemas de aeração. Pode recomendar
a instalação de viveiros automatizados. Orienta monitoramento constante
de parâmetros da água. Demonstra a forma de coletar amostras e de medir
temperatura, pH, oxigênio, transparência, entre outros parâmetros.
Indica o modo de instalar filtros de purificação e de dimensionar a
vazão da água. Monitora, antes do povoamento do viveiro com as
pós-lavas, a aplicação de calcário e de fertilizantes químicos e
orgânicos, que contribuirão na produção da alimentação natural, formada
por microorganismos, tais como algas, fito e zooplânctons, entre outros.
Além da alimentação natural, orienta o fornecimento de ração balanceada.
Seleciona e recomenda suplementos alimentares. Indica procedimentos
técnicos para aquisição de pós-larvas, para sua transferência ao local
de criação e para sua aclimatação. Pode organizar a reprodução de
camarões, monitorando a seleção, a mensuração e o acasalamento de
reprodutores. Pode orientar a aplicação de técnicas artificiais de
reprodução. Informa sobre técnicas de melhoramento genético. Recomenda o
emprego de tela de proteção durante o abastecimento de água no local da
criação, para impedir a entrada de espécies competidoras e predadoras.
Implanta controle da sanidade dos camarões, com observação de espécimes
e coleta de amostra dirigida - dos que apresentem aspectos menos
saudáveis - de diferentes pontos do viveiro. Orienta a pesagem dos
camarões e imediato envio para a avaliação de tecidos e órgãos, a fim de
identificar enfermidades virais e bacterianas, entre outras doenças e
irregularidades. Recomenda, ainda, o controle de índices zootécnicos -
como taxa de crescimento, peso corporal, entre outros -- indicativos da
saúde dos camarões em cativeiro. Orienta a despesca, a desinfecção e o
abate dos camarões. Organiza o beneficiamento dos camarões abatidos,
monitorando os processos, tais como limpeza, processamento em filés,
classificação, congelamento e embalagem. Monta a logística de
armazenagem e transporte da produção. Implanta e gerencia sistemas de
controle de qualidade de processos e produtos. Observa a organização do
local de trabalho. Monitora a aplicação de programas profiláticos,
higiênicos e sanitários. Orienta a conservação de ferramentas e
instrumentos em plenas condições de funcionamento. Prepara plano de
manutenção de máquinas e equipamentos. Elabora relatório técnico sobre
resultados do projeto de carcinicultura. Pode ministrar palestras. Pode
administrar empreendimentos e realizar pesquisa e extensão, na área de
carcinicultura. Pode supervisionar trabalho de equipe, avaliando
desempenho e desenvolvendo treinamentos. Mantém-se atualizado em relação
às inovações em carcinicultura. Estuda aplicação de novas tecnologias,
como o sistema de criação em bioflocos bacterianos para criação de
camarões marinhos em locais afastados da costa. Contribui para redução
dos impactos ambientais do processo de criação, monitorando a aplicação
das normas ambientais. Zela pela segurança, prevenindo situações de
risco. Utiliza e monitora o uso de equipamentos de proteção individual.
Presta primeiros socorros.

\begin{Form}
\textbf{Esta correspondência está adequada?} \\ 
\ChoiceMenu[radio,radiosymbol=\ding{52},name=cbovalid_ 321310 ]{}{CBO deve ficar=sim, \hspace{6em} CBO não fica aqui=nao}
\end{Form}

\textbf{CBO: 341220  ( Patrão de pesca de alto-mar )}

Realiza operações de navegação em mar aberto com embarcação de pesca,
traçando a rota de navegação; relatando posição da embarcação e hora
estimada de chegada; e monitorando condições de navegabilidade. Guarnece
o passadiço e o timão. Aciona equipamentos auxiliares e luzes
regulamentares da embarcação. Opera equipamentos de embarcação de pesca
na navegação em mar aberto, tais como equipamentos de comunicação, de
combate a incêndio, de salvatagem e de orientação para posicionamento
geográfico. Opera piloto automático, instrumento ecossondador e sonar.
Atraca e desatraca embarcação de pesca na navegação em mar aberto,
acionando a seção de máquinas; trocando informações com a estação de
apoio; e analisando as condições intervenientes (maré, vento, tráfego de
embarcações, visibilidade e profundidade). Solicita serviços de apoio
portuário. Realiza as manobras necessárias para atracar ou desatracar a
embarcação. Reboca embarcações. Monitora carga e descarga da embarcação
de pesca na navegação em mar aberto, emitindo notificação de prontidão
da embarcação; verificando condições dos porões de carga; classificando
cargas; e elaborando o plano de carregamento. Verifica documentação das
cargas. Designa pessoal para carregamento e descarregamento. Monitora o
calado, a disposição das cargas e as condições de conservação de
pescados. Verifica condições e fixação das cargas. Informa unidade
receptora sobre características das cargas. Supervisiona serviços de
manutenção de embarcação de pesca na navegação em mar aberto,
solicitando serviços de reparo e fiscalizando reparos realizados a
bordo. Fiscaliza condições de conservação da embarcação e condições de
equipamentos de convés. Fiscaliza equipamentos de combate a incêndio e
de salvatagem. Fiscaliza dotação de material para contenção de
poluentes. Fiscaliza compartimentos habitáveis. Administra recursos
materiais e financeiros nas operações com embarcação de pesca na
navegação em mar aberto, requisitando materiais, comprando insumos,
requisitando provisões de alimentos e requisitando combustível,
lubrificantes e água. Administra custeio de bordo. Inventaria materiais.
Efetua pagamento para a tripulação. Supervisiona tripulação de
embarcação de pesca na navegação em mar aberto, divulgando normas e
regulamentos, determinando horário de trabalho, distribuindo tarefas
para guarnição e monitorando atividades. Treina novatos. Aplica
advertências e demais punições. Comanda embarcação de pesca na navegação
em mar aberto, de acordo com normas da autoridade marítima da Diretoria
de Portos e Costas da Marinha do Brasil, que define capacidades das
embarcações e limites em que pode atuar. Registra dados da embarcação de
pesca na navegação em mar aberto, preenchendo rol de tripulação (rol de
equipagem); emitindo documentação de entrada, saída e permanência no
porto; e atualizando cartas e publicações náuticas. Organiza
documentação de embarcação e carga. Escritura diário de navegação.
Preenche livro de carga e documentação de despacho de lixo. Escritura
diário de comunicação e redige atas de reuniões. Preenche mapa de bordo
para o Instituto Brasileiro do Meio Ambiente e dos Recursos Naturais
Renováveis (Ibama). Orienta tripulação sobre temas de segurança do
trabalho, tais como utilização de Equipamentos de Proteção individual
(EPI); situações de emergência; e condições e atos inseguros. Simula
situações adversas para treinamento da tripulação. Divulga informações
sobre saúde e orienta sobre questões ambientais.

\begin{Form}
\textbf{Esta correspondência está adequada?} \\ 
\ChoiceMenu[radio,radiosymbol=\ding{52},name=cbovalid_ 341220 ]{}{CBO deve ficar=sim, \hspace{6em} CBO não fica aqui=nao}
\end{Form}

\newpage
\section{ 12001 -- Técnico em Defesa Civil }

\textbf{Perfil do Curso (CNCT)}

O Técnico em Defesa Civil será habilitado para:

Gerenciar riscos e desastres. Realizar o monitoramento preventivo de
desastres. Monitorar mudanças climáticas, alertas de emergências e
sistemas de informações geográficas. Coordenar, de forma estratégica,
secretarias, entidades e órgãos de poder público, privado e ONG.
Planejar reuniões de núcleos comunitários de proteção e defesa civil.
Orientar e mobilizar as comunidades a adotar comportamentos adequados de
prevenção, preparação, resposta e recuperação em situação de eventos
adversos/desastres e promover a autoproteção. Promover e coordenar ações
de recuperação de eventos adversos/desastres. Promover políticas
públicas para redução dos riscos de eventos adversos/desastres.
Estimular o desenvolvimento de cidades resilientes e os processos
sustentáveis de urbanização. Produzir alertas antecipados e possibilitar
mecanismos de comunicação com base no monitoramento sobre a
possibilidade de ocorrência de eventos adversos e/ou desastres naturais.
Oferecer capacitação de recursos humanos para as ações de proteção e
defesa civil. Elaborar planos de contingência de proteção e defesa
civil. Elaborar pareceres, relatórios, planos, projetos. Realizar
pesquisas, estudos, análises, interpretação, planejamento, implantação,
coordenação e controle dos trabalhos em proteção e Defesa Civil.

Para atuação como Técnico em Defesa Civil, são fundamentais:

Proficiência e conhecimento estratégico, tático e operacional.
Conhecimentos e saberes relacionados aos processos de projetos,
planejamento e gestão, tanto no setor público quanto no privado, em
situações de eventos adversos e possíveis desastres de origem natural ou
tecnológica. Conhecimento técnico para interpretar, monitorar e
gerenciar condições geológicas, meteorológicas, climatológicas,
epidemiológicas e outras condições fortuitas ligadas a condições de meio
ambiente, natureza e ambiente antropizados. Capacidade de assegurar a
saúde e a segurança dos cidadãos e a sustentabilidade do desenvolvimento
urbano. Conhecimentos e saberes relacionados às normas técnicas, à
liderança de equipes, à solução de problemas técnicos e trabalhistas e à
gestão de conflitos.

\textbf{CBOs não correspondidos e seus perfis:}

\textbf{CBO: 351605  ( Técnico em segurança do trabalho )}

Elabora, implanta e implementa política de saúde e segurança do
trabalho, definindo metas, prioridades e responsabilidades. Divulga a
política entre os trabalhadores. Acompanha e avalia os resultados,
podendo promover alterações para aperfeiçoamento e atualização. Analisa
e avalia o ambiente de trabalho, as instalações e os processos laborais,
identificando fatores de risco de acidentes do trabalho, de doenças
profissionais e de presença de agentes ambientais agressivos ao
trabalhador. Participa da seleção de tecnologias e processos de
trabalho, avaliando implantação e impactos por meio de análise coletiva
do trabalho (ACT) e análise ergonômica do trabalho (AET). Emite parecer
sobre equipamentos, máquinas e processos. Coleta, organiza e analisa
comunicações de acidentes de trabalho (CAT). Investiga e analisa
acidentes para determinar causas e recomendar medidas de prevenção e
controle. Levanta e utiliza dados estatísticos de doenças e acidentes de
trabalho para ajustes das ações prevencionistas. Adota medidas de
controle de riscos ocupacionais por meio de ações e programas de saúde e
segurança do trabalho. Realiza procedimentos de orientação sobre medidas
de eliminação e neutralização de riscos. Seleciona, controla, orienta e
fiscaliza o uso de EPC (Equipamento de Proteção Coletiva) e EPI
(Equipamento de Proteção Individual). Participa do desenvolvimento de
ações educativas e capacitações, elaborando cronograma, preparando
recursos didáticos e formando multiplicadores. Produz relatórios
referentes à segurança e à saúde do trabalhador. Elabora manual do
sistema de gestão de saúde e segurança no trabalho. Controla atualização
de documentos, normas e legislação. Atua em equipes multidisciplinares
de saúde e segurança ocupacionais.

\begin{Form}
\textbf{Esta correspondência está adequada?} \\ 
\ChoiceMenu[radio,radiosymbol=\ding{52},name=cbovalid_ 351605 ]{}{CBO deve ficar=sim, \hspace{6em} CBO não fica aqui=nao}
\end{Form}
\footnotetext{Fonte: CNCT e QBQ, 18/07/2025}

\end{document}
